% !TeX root = ../../Book2.tex
% 纲要标题：## 第8章　张量代数与对称代数

\chapter{张量代数与对称代数}
\label{b2:ch:08}
\BookChapterNumber{b2:ch:08}

\begin{chapterabstract}
多重张量积把多线性映射改写成线性映射，张量代数把所有次数的有序乘积放在同一分次代数中，对称代数则进一步施加交换关系。本章证明这些构造的泛性质，建立自由张量积的基与有限自由情形的多项式解释，并比较对称商与不变张量，说明阶乘可逆时的同构及这一条件缺失时可能出现的差别。最后把对偶、张量积与对称幂用于群表示，说明等变映射如何成为不变张量，以对偶空间的对称代数定义多项式不变环，并比较位置置换与一般线性群作用的对称算子。
\end{chapterabstract}

\begin{learninggoals}
\begin{itemize}
\item 用多重张量积表示多线性映射，说明改变括号与置换因子的指定同构及其相容性。
\item 构造分次张量代数，区分一般模的张量代数与选定自由生成元后的字代数，利用生成元和关系定义同态。
\item 证明对称代数的泛性质及多项式解释，计算对称幂的基和维数，并检查阶乘是否可逆。
\item 写出张量、对偶、对称幂表示的作用公式，区分位置置换不变与群作用不变，将有限维 Hom 识别为张量表示。
\end{itemize}
\end{learninggoals}

取自由模 \(V=Re_1\oplus Re_2\)。在张量代数中，\(e_1e_2\) 与 \(e_2e_1\) 记录两种不同次序；若希望它们相同，就必须再施加交换关系，得到对称代数中的同一个乘积。这不是改变记号而已：给两个元素指定像时，映到一般结合代数无需它们可交换，映到交换代数才自动满足新关系。不同的关系决定了构造可以接受哪些映射。

\ProseParagraphBreak

这里还要区分两种构造：一种选出在位置置换下已经固定的张量，另一种把不同排列识别为同一个类。它们将在\secref{b2:sec:0803}中分别定义；次数二的计算先说明为什么不能无条件混同。即使取 \(R=\mathbb Z\)，张量 \(e_1\otimes e_2+e_2\otimes e_1\) 在交换商中的像仍是 \(2e_1e_2\)，不能由它得到 \(e_1e_2\)。所以在二不可逆时，不能只凭对称性就断言不变张量与对称商相同；在一般次数中，阶乘可逆将提供统一的平均方法；不可逆时，则须直接检查从固定子模到商的映射。随后还要区分位置置换与群在各因子上同时作用，这两种作用所固定的张量也未必相同。

\ProseParagraphBreak

本章以第六章的张量积泛性质、剩余作用和结合同构为直接基础；自然性采用第五章的定义，有限自由模的坐标与对偶采用第一、三章的结果。除表示部分把底环取为域外，前面大部分构造只需交换底环。各类商构造的泛性质都要同时说明所除去的关系，以及从商出发的映射为何零化这些关系。

\ProseParagraphBreak
张量代数与对称代数的商构造可参阅 Etingof 等人的《Introduction to Representation Theory》\cite[Section 2.12, Definition 2.12.1, p. 35]{EtingofEtAl2011RepresentationTheory}；对偶与张量积表示的公式见\cite[Section 4.4, p. 65]{EtingofEtAl2011RepresentationTheory}。对称幂的商定义适用于一般交换环；用平均算子识别不变张量时，相关阶乘须可逆。

% !TeX root = ../../Book2.tex
% 纲要标题：### 8.1 多线性映射

\section{多线性映射}
\label{b2:sec:0801}

从双线性映射走向多线性映射，首先遇到的不是新运算，而是括号和标量的安排。一个三变量公式可以先合并前两个变量，也可以先合并后两个变量；只有证明两种构造通过指定同构相容，才可以省略括号。本节固定交换环 \(R\)，所有模均为幺性 \(R\) 模。这一交换性约定允许在不同因子间移动标量；第六章一般双模的作用侧别仍是它的基础。

% 纲要要点（供目录核对）：
%
% * 多重张量积
% * 自然同构与结合约定
% * 自由模张量积的基与有限秩情形的秩公式
%

\subsection{多重张量积}
\label{b2:subsec:0801:01}

\begin{definition}\label{b2:def:multilinear-map}
设 \(R\) 为交换环，\(M_1,\ldots,M_d,N\) 为 \(R\) 模，\(d\ge1\)。如果映射
\[
b:M_1\times\cdots\times M_d\longrightarrow N
\]
在固定其余变量后对每个变量都是 \(R\) 线性的，就称 \(b\) 为 \(R\) 上的\term[duoxianxing-yingshe]{多线性映射}[multilinear map]。具体说，对每个 \(1\le j\le d\)，任取 \(r,s\in R\)、\(x,y\in M_j\) 及其余各位置的 \(m_i\in M_i\)，都有
\[
\begin{aligned}
b(m_1,\ldots,rx+sy,\ldots,m_d)
&=r\cdot b(m_1,\ldots,x,\ldots,m_d)\\
&\quad+s\cdot b(m_1,\ldots,y,\ldots,m_d).
\end{aligned}
\]
\end{definition}

直积 \(M_1\times\cdots\times M_d\) 本身是逐坐标运算的 \(R\) 模。当 \(d=1\) 时，多线性就是线性；当 \(d\ge2\) 时，对每个变量分别线性与对这个直积模整体线性是不同的要求。多线性映射满足 \(b(rm_1,\ldots,rm_d)=r^d\cdot b(m_1,\ldots,m_d)\)，整体线性却要求右端只有一个标量 \(r\)。例如 \(R=\mathbb R\)，乘积 \((x,y,z)\mapsto xyz\) 对每个变量分别线性，但把三个变量同时乘二会使值乘八，而非乘二。

\ProseParagraphBreak
令 \(T_1=M_1\)，递归定义 \(T_d=T_{d-1}\otimes_RM_d\)。因 \(R\) 交换，第六章的剩余作用使每个 \(T_j\) 为 \(R\) 模，因而可以继续取张量积。把 \((m_1,\ldots,m_d)\) 映到
\[
\Bigl(\cdots\bigl((m_1\otimes m_2)\otimes m_3\bigr)\cdots\Bigr)\otimes m_d
\]
的映射记作 \(\tau\)。以后用 \(m_1\otimes\cdots\otimes m_d\) 表示这个元素，用 \(\bigotimes_{j=1}^d M_j\) 表示 \(T_d\)。左结合的递归构造先给出了一个具体模型；下面的多线性泛性质将刻画这个对象，使后续构造不再依赖选用哪一种括号形式。

\begin{theorem}[多重张量积的泛性质]\label{b2:thm:multiple-tensor-universal}
设 \(R\) 为交换环，\(d\ge1\)，\(M_1,\ldots,M_d\) 为 \(R\) 模。令 \(\tau:M_1\times\cdots\times M_d\to\bigotimes_{j=1}^dM_j\) 把 \((m_1,\ldots,m_d)\) 送到纯张量 \(m_1\otimes\cdots\otimes m_d\)，则 \(\tau\) 是多线性映射。对任意 \(R\) 模 \(N\) 与多线性映射 \(b:M_1\times\cdots\times M_d\rightarrow N\)，存在唯一 \(R\) 线性映射
\[
\widetilde b:\bigotimes_{j=1}^dM_j\longrightarrow N,
\qquad \widetilde b(m_1\otimes\cdots\otimes m_d)=b(m_1,\ldots,m_d).
\]
多重张量积由这些纯张量生成，并在保持指定多线性映射的唯一同构意义下唯一。
\end{theorem}
\begin{proof}
三变量时，给定三线性映射 \(b:M_1\times M_2\times M_3\rightarrow N\)，固定 \(z\in M_3\) 后，二重张量积的泛性质给出唯一线性映射
\[
b_z:M_1\otimes_RM_2\longrightarrow N,
\qquad b_z(x\otimes y)=b(x,y,z).
\]
对 \(r,s\in R\) 及 \(z,w\in M_3\)，映射 \(b_{rz+sw}\) 与 \(r b_z+s b_w\) 在每个 \(x\otimes y\) 上取值相同，故由纯张量生成性相等。因此 \((t,z)\mapsto b_z(t)\) 是 \((M_1\otimes_RM_2)\times M_3\) 上的双线性映射，再应用一次张量积的泛性质，就得到 \((M_1\otimes_RM_2)\otimes_RM_3\rightarrow N\)。这样先合并前两个变量，再合并最后一个变量，并没有改变原来的三线性取值。

对一般 \(d\) 按同一方法归纳。\(d=1\) 时是恒等映射，\(d=2\) 时是\cref{b2:thm:tensor-universal}及\cref{b2:prop:tensor-bimodule}给出的交换环线性版本。设结论对 \(d-1\) 成立，固定 \(m_d\) 后得到唯一线性映射 \(b_{m_d}:T_{d-1}\rightarrow N\)。这些映射随 \(m_d\) 线性变化：所需等式在前 \(d-1\) 重纯张量上由 \(b\) 的多线性成立，再由生成性推广到整个 \(T_{d-1}\)。故 \((t,m_d)\mapsto b_{m_d}(t)\) 为双线性映射。

应用\cref{b2:thm:tensor-universal}及\cref{b2:prop:tensor-bimodule}给出的交换环线性版本，得到所需 \(\widetilde b\)。任一 \(T_d\) 中的元素先写成有限和 \(\sum t_i\otimes m_i\)，再把各 \(t_i\) 展开成前 \(d-1\) 重纯张量，故 \(d\) 重纯张量生成全模。这证明映射唯一，也说明 \(\tau\) 的逐变量线性由每一步的双线性继承。

若另一个对象带有相同泛性质，把两者的指定多线性映射分别代入对方的泛性质，得到正反两个线性映射。两复合在每个指定纯张量上为恒等，故在全模为恒等；保持指定映射的同构也由生成性唯一确定。
\end{proof}

为了统一后面的次数与乘法，采用零重张量积为 \(R\) 的约定，并把没有输入变量的“多线性映射”约定为指定目标模中的一个元素 \(n\)。它对应唯一线性映射 \(R\rightarrow N\)、\(r\mapsto rn\)，其在 \(1\) 处的值为 \(n\)。次数零的这一约定将提供张量代数的单位元。

\subsection{自然同构与结合约定}
\label{b2:subsec:0801:02}

任意完整括号形式都可递归应用二重张量积构造。对因子个数归纳，把最外层写成 \(X\otimes_RY\)，其中 \(X,Y\) 是较少因子的括号积。固定后半组变量，由 \(X\) 的泛性质合并前半组；所得映射对后半组仍逐变量线性，再由 \(Y\) 的泛性质合并，得到 \(X\times Y\) 上的双线性映射。二重张量积的泛性质给出所需线性映射，唯一性则由有序纯张量生成全模得到。因此每种括号形式都满足\cref{b2:thm:multiple-tensor-universal}的同一个 \(d\) 变量泛性质，它们之间有唯一同构，将相同顺序的纯张量送成彼此。\cref{b2:thm:tensor-associativity-unit}的结合同构正是三因子时的这个映射。

\ProseParagraphBreak
起点与终点的括号形式给定后，任何逐步改变括号的路径都保持同一组有序纯张量。由泛性质的唯一性，这些复合映射必定相同；第六章的四因子计算就是这一事实的具体情形。因此下文可以省略括号：无括号的记号表示经这些典范同构识别的张量积，而非宣称不同递归模型字面相等。

\begin{proposition}\label{b2:prop:multiple-tensor-functoriality}
设 \(R\) 为交换环，\(d\ge1\)，\(M_j,N_j\) 为 \(R\) 模。给定 \(R\) 线性映射 \(f_j:M_j\rightarrow N_j\)（\(1\le j\le d\)），存在唯一 \(R\) 线性映射
\[
\bigotimes_{j=1}^d f_j:\bigotimes_{j=1}^d M_j\longrightarrow\bigotimes_{j=1}^d N_j,
\qquad \bigotimes_jm_j\longmapsto\bigotimes_j f_j(m_j).
\]
它保持恒等映射与逐因子复合，并与保持有序纯张量的结合同构相容。若 \(\sigma\in S_d\)，还存在置换因子的同构
\[
P_\sigma:M_1\otimes\cdots\otimes M_d
\longrightarrow M_{\sigma^{-1}(1)}\otimes\cdots\otimes M_{\sigma^{-1}(d)},
\quad \bigotimes_jm_j\longmapsto
m_{\sigma^{-1}(1)}\otimes\cdots\otimes m_{\sigma^{-1}(d)}.
\]
当所有因子相同时，有 \(P_\sigma P_\pi=P_{\sigma\pi}\)。
\end{proposition}
\begin{proof}
所给公式在原变量中多线性，所以由多重张量积的泛性质诱导线性映射。对逐因子复合，两个可能映射都把 \(\bigotimes_jm_j\) 送到 \(\bigotimes_j g_j\bigl(f_j(m_j)\bigr)\)，故相等；恒等映射与结合的相容性由相同的纯张量检验得到。

置换公式仍逐变量线性，因此也给出映射，反向用 \(\sigma^{-1}\) 即得逆。所有因子相同时，连续做 \(P_\pi\)、\(P_\sigma\)，第 \(j\) 个位置最终来自原来的 \(\pi^{-1}\sigma^{-1}(j)=(\sigma\pi)^{-1}(j)\) 位置，所以满足所述乘法公式。
\end{proof}

\symidx{\(P_\sigma\)：置换张量因子的算子}
逆置换使原来第 \(j\) 个位置的元素移动到第 \(\sigma(j)\) 个位置，从而给出左群作用。例如在 \(M^{\otimes3}\) 中，取 \(\sigma=(123)\)、\(\pi=(12)\)，并按先做 \(\pi\) 再做 \(\sigma\) 的次序复合，则
\[
\begin{aligned}
P_\sigma(a\otimes b\otimes c)&=c\otimes a\otimes b,\\
P_\sigma P_\pi(a\otimes b\otimes c)
&=P_\sigma(b\otimes a\otimes c)
=c\otimes b\otimes a.
\end{aligned}
\]
这正是 \(P_{\sigma\pi}=P_{(13)}\) 的取值。若把公式中的 \(\sigma^{-1}\) 换成 \(\sigma\)，复合反而会对应 \(\pi\sigma\)，次序被倒转。置换因子也没有交换代数中两个元素的含义：\(M\otimes N\cong N\otimes M\) 是指定同构，通常不能在一个尚未给乘法的模中把 \(m,n\) 相乘。

\subsection{有限自由情形的基与维数}
\label{b2:subsec:0801:03}

纯基张量的基定理不要求各因子的基有限；有限性只在把所得基的大小写成有限秩或维数公式时需要。

\begin{proposition}\label{b2:prop:multiple-tensor-basis}
设 \(R\) 为交换环，\(d\ge1\)，各 \(M_j\) 为自由 \(R\) 模，具有指定基 \((e_{j,i})_{i\in I_j}\)，则
\[
e_{1,i_1}\otimes\cdots\otimes e_{d,i_d},
\qquad (i_1,\ldots,i_d)\in I_1\times\cdots\times I_d,
\]
构成多重张量积的一组基。若 \(R\ne0\)，且各 \(I_j\) 有限，记 \(n_j=\lvert I_j\rvert\)，则其秩为 \(n_1\cdots n_d\)；在域上，这就是有限维空间的维数乘法公式。
\end{proposition}
\begin{proof}
逐因子展开 \(m_j=\sum_{i_j\in I_j}a_{j,i_j}e_{j,i_j}\)。每个展开只有有限多个非零系数，而因子的个数 \(d\) 有限，所以多线性给出有限和
\[
m_1\otimes\cdots\otimes m_d
=\sum_{i_1,\ldots,i_d}
a_{1,i_1}\cdots a_{d,i_d}\,
e_{1,i_1}\otimes\cdots\otimes e_{d,i_d}.
\]
纯张量生成全模，故所列元素生成。对每组下标 \((i_1,\ldots,i_d)\)，将各因子的第 \(i_j\) 个坐标相乘，得到多线性映射 \((m_1,\ldots,m_d)\mapsto a_{1,i_1}\cdots a_{d,i_d}\)。它诱导的线性泛函在对应的纯基张量上取一，在其他纯基张量上取零。对任意有限线性关系应用相应的坐标泛函，逐项得到所有系数为零，故线性无关。基有限时共有 \(n_1\cdots n_d\) 个纯基张量；秩的含义由\cref{b2:thm:free-rank-invariance}保证。
\end{proof}

例如 \(R^2\otimes_RR^3\otimes_RR^2\) 有十二个标准纯基张量。在 \(R=\mathbb Z\) 时，
\[
(e_1+2e_2)\otimes(f_1-f_3)\otimes e_2
=e_1\otimes f_1\otimes e_2-e_1\otimes f_3\otimes e_2
+2e_2\otimes f_1\otimes e_2-2e_2\otimes f_3\otimes e_2.
\]
这些系数是唯一的，但将一个张量写成几个任意纯张量之和并不唯一。\subsecref{b2:subsec:0601:03}中已经证明，\(k^2\otimes_k k^2\) 中的 \(e_1\otimes e_1+e_2\otimes e_2\) 不能写成一个纯张量。因此，基张量坐标和一般纯张量分解是两种不同的数据。

\ProseParagraphBreak
自由假设不能只替换成“有限生成”。例如 \(M=\mathbb Z/2\mathbb Z\)、\(N=\mathbb Z/3\mathbb Z\) 均可由一个元素生成，\cref{b2:prop:tensor-quotient}却给出 \(M\otimes_{\mathbb Z}N=0\)。没有基时，张量关系可能消去整个模，不能用生成元数相乘代替秩或维数。

% !TeX root = ../../Book2.tex
% 纲要标题：### 8.2 张量代数

\section{张量代数}
\label{b2:sec:0802}

固定一个模 \(M\) 后，可以把不同次数的张量同时放进一个对象：一次张量是 \(M\) 的元素，二次张量记录两个元素的有序相乘，三次张量记录三个元素的有序相乘。把这些次数作直和，再用连接张量串定义乘法，便得到张量代数。这里首先保留次序；交换性将作为下一节另外施加的关系。

% 纲要要点（供目录核对）：
%
% * 分次张量代数
% * 自由结合代数的泛性质
% * 商代数与按关系定义代数
% * 增广理想的一次商与底模恢复
%

\subsection{分次张量代数}
\label{b2:subsec:0802:01}

\begin{definition}\label{b2:def:associative-graded-algebra}
设 \(R\) 为交换环。给定含单位元的环 \(A\) 及保单位元环同态 \(R\rightarrow Z(A)\)，称 \(A\) 为 \(R\) 上的\term[jiehe-daishu]{结合代数}[associative algebra]。它因此是 \(R\) 模，乘法对两个变量分别 \(R\) 线性；结构同态不要求单射。两个 \(R\) 代数之间与各自结构同态相容的保单位元环同态，称为 \(R\) 代数同态。

若另有 \(R\) 模直和分解
\[
A=\bigoplus_{d\ge0}A_d,\qquad 1_A\in A_0,\qquad A_iA_j\subseteq A_{i+j},
\]
则称 \(A\) 为非负整数\term[fenci-daishu]{分次代数}[graded algebra]；\(A_d\) 中的元素称为次数 \(d\) 的\term[qici-yuansu]{齐次元素}[homogeneous element]。保持各次数分量的代数同态称为分次同态。
\end{definition}

结构同态给出标量的作用，并不表示 \(A\) 必须含有一个与 \(R\) 相同的子环。例如对 \(n>1\)，商映射 \(\mathbb Z\rightarrow\mathbb Z/n\mathbb Z\) 使后者成为 \(\mathbb Z\) 代数，而有限环 \(\mathbb Z/n\mathbb Z\) 不可能含有 \(\mathbb Z\) 子环。

\ProseParagraphBreak
直和要求一个元素只有有限多个非零齐次分量，因此不包含任意无穷和。分次也没有规定 \(ab=(-1)^{ij}ba\)；这种带符号的交换条件是另一个额外性质。普通交换的多项式代数以全次数分次，张量代数则一般不交换。

\begin{definition}\label{b2:def:tensor-algebra}
设 \(R\) 为交换环，\(M\) 为 \(R\) 模。令 \(M^{\otimes0}=R\)、\(M^{\otimes1}=M\)，并记
\[
T_R(M)=\bigoplus_{d\ge0}M^{\otimes d}.
\]
对正次数的纯张量规定
\[
(m_1\otimes\cdots\otimes m_i)(n_1\otimes\cdots\otimes n_j)
=m_1\otimes\cdots\otimes m_i\otimes n_1\otimes\cdots\otimes n_j,
\]
对 \(r\in R=M^{\otimes0}\) 与 \(x\in M^{\otimes d}\)（\(d\ge0\)），规定 \(rx=xr=r\cdot x\)，再按双线性扩张到整个直和。所得分次 \(R\) 代数称为 \(M\) 的\term[zhangliang-daishu]{张量代数}[tensor algebra]。
\end{definition}
\symidx{\(M^{\otimes d}\)：模的 \(d\) 次张量幂}
\symidx{\(T_R(M)\)：模的张量代数}

各个 \(M^{\otimes d}\) 原本只是不同次数的模；上述连接把两个张量串变成一个更长的串，才给这些次数之间定义了乘法。其良定义由张量积的泛性质保证。对每对 \(i,j>0\)，\cref{b2:thm:multiple-tensor-universal}与结合同构给出线性映射
\[
M^{\otimes i}\otimes_RM^{\otimes j}\longrightarrow M^{\otimes(i+j)}
\]
并给出上面的纯张量公式，所以它不依赖纯张量和的表达。任意两个元素仅有有限多个分量，将这些双线性乘法逐项相加仍为有限和。三个齐次元素相乘时，两种结合方式都连接同一个有序张量串；由纯张量生成性及分配律，结合律在所有元素上成立。\(1_R\) 是单位，标量与所有张量可交换，次数相加，因此定义确实给出分次 \(R\) 代数。

\ProseParagraphBreak
一次分量的包含记为 \(\iota:M\rightarrow T_R(M)\)，它单射，因为这是一个直和分量的包含。\(T_R(M)\) 由 \(\iota(M)\) 和 \(R\) 作为代数生成：每个纯张量都等于若干一次元素之积，任意元素又是这些纯张量与常数的有限线性组合。

\subsection{自由结合代数的泛性质}
\label{b2:subsec:0802:02}

\begin{theorem}[张量代数的泛性质]\label{b2:thm:tensor-algebra-universal}
设 \(R\) 为交换环，\(M\) 为 \(R\) 模，\(\iota:M\hookrightarrow T_R(M)\) 为张量代数的一次分量包含。对任意含单位元的结合 \(R\) 代数 \(A\) 及 \(R\) 线性映射 \(f:M\rightarrow A\)，存在唯一 \(R\) 代数同态
\[
\widehat f:T_R(M)\longrightarrow A,\qquad \widehat f\iota=f.
\]
其零次部分为 \(A\) 的结构同态，其正次部分满足
\[
\widehat f(m_1\otimes\cdots\otimes m_d)=f(m_1)\cdots f(m_d).
\]
\end{theorem}
\begin{proof}
由于 \(R\) 的像在 \(A\) 中心，右端对 \(m_1,\ldots,m_d\) 的每个变量都 \(R\) 线性。\cref{b2:thm:multiple-tensor-universal}给出唯一线性映射 \(f_d:M^{\otimes d}\rightarrow A\)。令 \(f_0\) 为结构同态，利用直和的泛性质把它们合成 \(\widehat f\)。在两个纯张量上，先连接再作用与分别作用后在 \(A\) 相乘相同；由分配律推广到有限和，故 \(\widehat f\) 保持乘法。它在零次部分保持单位元及标量，所以是 \(R\) 代数同态。

任何所求同态都必须在一次分量上等于 \(f\)，在纯张量上等于这些像的有序乘积，在 \(R\) 上等于结构同态。因此它在生成全代数的元素上已被确定，唯一性成立。
\end{proof}

该泛性质的目标不要求交换，也不要求 \(f\) 单射。张量代数的“自由”是说：\(M\) 中原有的线性关系仍然保留，但不再对它的元素另加乘法关系。因此任意到结合 \(R\) 代数的线性映射都能唯一延拓，而不需要额外检查这些像之间满足何种乘法关系。

\ProseParagraphBreak
将一次元素全部送到零，就得到\term[zengguang-tongtai]{增广同态}[augmentation] \(\varepsilon:T_R(M)\rightarrow R\)。它是取零次分量的投影：若 \(x_d\in M^{\otimes d}\)，则
\[
\varepsilon(r+x_1+\cdots+x_n)=r.
\]
其核为正次部分 \(I=T_R(M)_+=\bigoplus_{d\ge1}M^{\otimes d}\)，称为增广理想。\cref{b2:prop:tensor-augmentation-indecomposables}将说明，模掉其中至少含两个一次因子的乘积后，\(I/I^2\) 恰好恢复一次分量 \(M\)。

\ProseParagraphBreak
把 \(R\) 线性映射 \(f:M\rightarrow N\) 后接 \(N\) 到其张量代数的包含，得到分次同态 \(T_R(f):T_R(M)\rightarrow T_R(N)\)，它在 \(d\) 次分量上就是 \(f^{\otimes d}\)，在零次分量上为 \(R\) 的恒等映射。恒等及复合在一次生成元上核对，便知 \(T_R\) 是函子。

\ProseParagraphBreak
这个函子却不一定保持单射。第七章用来检验平坦性的双数环（\cref{b2:prop:dual-number-flat}）给出一个具体例子：设 \(k\) 为域，\(R=k[\delta]/(\delta^2)\)，取 \(M=(\delta)\) 及包含 \(i:M\hookrightarrow R\)。映射 \(R\rightarrow M\)、\(a\mapsto a\delta\) 的核为 \((\delta)\)，故 \(M\cong R/(\delta)\)。由\cref{b2:prop:tensor-quotient}，
\[
M\otimes_RM\cong R/(\delta),
\]
其中 \(\delta\otimes\delta\) 对应 \(\bar1\)，所以非零。它经 \(i\otimes i\) 映到 \(R\otimes_RR\) 后却变成
\[
\delta\otimes\delta=1\otimes\delta^2=0.
\]
因此 \(T_R(i)\) 已在二次分量中有非零核。泛性质保证线性映射可以延拓成代数同态，不能保证单射也延拓成单射。

\begin{proposition}\label{b2:prop:free-associative-algebra}
设 \(R\) 为交换环，\(M\) 为以 \((x_i)_{i\in I}\) 为基的自由 \(R\) 模，则 \(T_R(M)\) 以所有有限字
\[
1,\quad x_{i_1},\quad x_{i_1}x_{i_2},\quad\ldots
\]
为 \(R\) 模的一组基，乘法为连接字。给定任意 \(R\) 代数 \(A\) 中的一族元素 \((a_i)_{i\in I}\)，存在唯一 \(R\) 代数同态将每个 \(x_i\) 送到 \(a_i\)。因此称 \(T_R(M)\) 为生成元 \((x_i)_{i\in I}\) 上的\term[ziyou-jiehe-daishu]{自由结合代数}[free associative algebra]，记为 \(R\langle x_i\mid i\in I\rangle\)。
\end{proposition}
\begin{proof}
由适用于任意自由基的\cref{b2:prop:multiple-tensor-basis}，每个正次 \(M^{\otimes d}\) 以长度为 \(d\) 的字为基，零次部分 \(R\) 则以空字 \(1\) 为基。不同次数直和后，全部有限字构成基。任意指定 \(x_i\mapsto a_i\) 由自由模的泛性质唯一延拓为 \(M\rightarrow A\)，再由\cref{b2:thm:tensor-algebra-universal}唯一延拓为代数同态。
\end{proof}
\symidx{\(R\langle x_i\mid i\in I\rangle\)：非交换生成元上的自由结合代数}

当 \(R\ne0\) 且有两个不同基向量 \(x,y\) 时，\(xy\)、\(yx\) 是不同的基字，所以 \(xy-yx\ne0\)。只有一个生成元时，每个字由长度决定，\(R\langle x\rangle\cong R[x]\)。一般的 \(T_R(M)\) 自由地接受来自模 \(M\) 的线性映射；只有 \(M\) 是以 \((x_i)\) 为基的自由模时，才可以把它进一步看成符号集合 \(\{x_i\}\) 上不带线性关系的字代数。非自由 \(M\) 中已有的线性关系，在张量代数的一次分量中仍然成立。

\subsection{商代数与按关系定义代数}
\label{b2:subsec:0802:03}

设 \(A\) 是 \(R\) 代数，\(J\) 是双边理想。因为结构标量在中心，\(J\) 自动为 \(R\) 子模，商环 \(A/J\) 继承 \(R\) 代数结构。第三章的模展示只需把关系的线性组合置为零；代数展示还必须让关系在左右乘法下继续成立：若 \(r_\alpha=0\)，则对任意 \(a,b\in A\)，也必须有 \(ar_\alpha b=0\)。所以给定自由代数中的关系族 \((r_\alpha)\) 后，必须取其生成的双边理想 \((r_\alpha)\)，它由有限和 \(\sum_\nu a_\nu r_{\alpha_\nu}b_\nu\) 组成。商代数
\[
R\langle x_i\mid i\in I\rangle/(r_\alpha)
\]
就是按生成元与关系给出的\term[daishu-zhanshi]{代数展示}[presentation of an algebra]，也称代数呈示。\termidx[daishu-chengshi]{代数呈示}这里的关系是代数乘法关系，也可含常数项，不能只把它们看成一次模关系。

\begin{proposition}\label{b2:prop:algebra-relations-universal}
设 \(R\) 为交换环，\(I\) 为指标集，\((r_\alpha)\) 是自由结合 \(R\) 代数 \(R\langle x_i\mid i\in I\rangle\) 中的一族元素，并以 \((r_\alpha)\) 记其生成的双边理想。对任意 \(R\) 代数 \(B\)，从 \(R\langle x_i\mid i\in I\rangle/(r_\alpha)\) 到 \(B\) 的 \(R\) 代数同态，与满足全部关系 \(r_\alpha(b_i)=0\) 的元素族 \((b_i)_{i\in I}\subseteq B\) 一一对应。
\end{proposition}
\begin{proof}
记 \(J=(r_\alpha)\) 为陈述中的双边理想。由\cref{b2:prop:free-associative-algebra}，给定像族恰好确定自由代数上的同态 \(h\)。若全部 \(h(r_\alpha)=0\)，则
\[
h\left(\sum_\nu a_\nu r_{\alpha_\nu}b_\nu\right)
=\sum_\nu h(a_\nu)h(r_{\alpha_\nu})h(b_\nu)=0,
\]
故它将关系理想送到零。公式 \(\bar h(a+J)=h(a)\) 因而不依赖代表元，给出唯一商代数同态。反向把商同态与商映射复合，其像族必满足全部关系。两个过程互逆。
\end{proof}

因此，从展示代数定义同态时，给定生成元的像就确定了唯一性，检查这些像满足全部关系则保证存在性。若还希望商代数保留分次，就必须保证关系理想包含其每个元素的全部齐次分量；否则一条混合不同次数的关系可能把原本应当分开的分量联系起来。由齐次关系生成理想可以保证这一点。

\begin{proposition}\label{b2:prop:homogeneous-ideal-quotient}
设 \(R\) 为交换环，\(A=\bigoplus_{d\ge0}A_d\) 为非负整数分次 \(R\) 代数。由 \(A\) 中齐次元素族生成的双边理想 \(J\) 满足
\[
J=\bigoplus_{d\ge0}(J\cap A_d).
\]
因此商代数有分次 \(A/J\cong\bigoplus_{d\ge0}A_d/(J\cap A_d)\)，商映射保持次数。
\end{proposition}
\begin{proof}
设生成关系 \(r_\alpha\) 为齐次。\(J\) 中元素是 \(a r_\alpha b\) 的有限和；把 \(a,b\) 各自分解成有限个齐次分量，每个乘积都为某个次数的齐次元素，且仍在 \(J\)。所以一个 \(J\) 元素的每个齐次分量都属于 \(J\)，得到所述直和等式。

把各 \(A_d\) 的商映到 \(A/J\) 后，每个商元素都有有限齐次表达。若这样的和为零，则其原像之和在 \(J\)，刚证的分量性质使每个原像分量都在 \(J\cap A_d\)，故各商分量为零。乘法仍按次数相加，由此得到所述分次。
\end{proof}

这样的理想称为\term[qici-lixiang]{齐次理想}[homogeneous ideal]。例如关系 \(xy-yx\) 齐次，次数为二。关系 \(yx-xy-1\) 却混合次数二与零，商后不能继续把 \(x,y\) 都放在次数一，除非单位元也变成了零。为排除这种零环坍缩，下面只需构造一个满足关系且单位元非零的代数。在任意域 \(k\) 上，令 \(X:k[t]\rightarrow k[t]\) 为乘 \(t\)，\(D\) 为形式求导，则对多项式 \(f\)，
\[
(DX-XD)(f)=(tf)'-tf'=f.
\]
所以由\cref{b2:prop:algebra-relations-universal}，赋值 \(x\mapsto X\)、\(y\mapsto D\) 给出 \(k\langle x,y\rangle/(yx-xy-1)\) 到 \(\operatorname{End}_k(k[t])\) 的代数同态，将单位元映成非零恒等算子。于是商中的单位元确实非零；若这个商继承了原次数分次，等式 \(yx-xy=1\) 将使一个非零元素同时属于次数二和零，与分次直和矛盾。微分算子在这里的作用就是证明商非零，从而确认不同次数相混合的障碍实际存在。

\ProseParagraphBreak
在张量代数中，增广理想记录所有正次部分，而它的平方只保留至少含两个一次因子的部分。因此取一次商可以在 \(R\) 模同构的意义下，从带有增广同态的代数恢复原来的模。

\begin{proposition}\label{b2:prop:tensor-augmentation-indecomposables}
设 \(R\) 为交换环，\(M\) 为 \(R\) 模，\(T_R(M)=\bigoplus_{d\ge0}M^{\otimes d}\) 为张量代数，\(\iota:M\hookrightarrow T_R(M)\) 为一次分量包含。由零映射 \(M\rightarrow R\) 诱导的增广同态 \(\varepsilon:T_R(M)\rightarrow R\) 是零次分量的投影。其核 \(I\) 满足
\[
I=\bigoplus_{d\ge1}M^{\otimes d},
\qquad I^2=\bigoplus_{d\ge2}M^{\otimes d},
\]
且映射
\[
M\longrightarrow I/I^2,
\qquad m\longmapsto\iota(m)+I^2
\]
是 \(R\) 模同构，且对 \(M\) 的 \(R\) 线性映射自然。
\end{proposition}
\begin{proof}
取零次分量的投影 \(p:T_R(M)\rightarrow R\) 保持单位元及标量。两个元素相乘时，由非负次数相加，乘积的零次分量恰为两者零次分量之积，所以 \(p\) 是 \(R\) 代数同态。它在一次分量上为零，故由\cref{b2:thm:tensor-algebra-universal}的唯一性等于 \(\varepsilon\)，并给出所述核 \(I\)。

任意两个 \(I\) 元素相乘，所有分量的次数至少为二，所以 \(I^2\subseteq\bigoplus_{d\ge2}M^{\otimes d}\)。反过来，对 \(d\ge2\)，每个纯张量都能写成
\[
m_1\otimes\cdots\otimes m_d
=\iota(m_1)(m_2\otimes\cdots\otimes m_d),
\]
右端两个因子均属于 \(I\)。由纯张量生成性，整个 \(M^{\otimes d}\) 都在 \(I^2\) 中，故得反向包含。于是 \(I=\iota(M)\oplus I^2\) 作为 \(R\) 模成立，取商即得所述同构。这个证明不要求 \(M\) 自由。

给定 \(R\) 线性映射 \(f:M\rightarrow N\)，记 \(I'\) 为 \(T_R(N)\) 的增广理想，\(\iota_N\) 为其一次分量包含。分次同态 \(T_R(f)\) 将 \(I\) 映入 \(I'\)，并将 \(I^2\) 映入 \((I')^2\)；它在一次商上把 \(\iota(m)+I^2\) 送到 \(\iota_N(f(m))+(I')^2\)。这与先作用 \(f\) 再取上述同构相同，故同构自然。
\end{proof}

% !TeX root = ../../Book2.tex
% 纲要标题：### 8.3 对称代数

\section{对称代数}
\label{b2:sec:0803}

张量代数保留因子的排列顺序。若把 \(uv\) 与 \(vu\) 识别，就得到只记录交换乘积的代数。这个商构造在任意交换底环上都成立；把它认成张量幂中的“不变部分”，却通常需要能除以阶乘。本节把这两个对象分别定义，再说明何时存在指定的同构。

% 纲要要点（供目录核对）：
%
% * 对称代数的泛性质
% * 有限自由模上的多项式解释
% * 对称幂的余不变量解释、对称化及特征 p 的取幂
% * 三次对称化的轨道计算及特征二、三的比较
% * 张量代数与对称代数的扩张标量相容性
%

\subsection{对称代数的泛性质}
\label{b2:subsec:0803:01}

\begin{definition}\label{b2:def:symmetric-algebra}
设 \(R\) 为交换含幺环，\(M\) 为幺性 \(R\) 模。令 \(J\) 为 \(T_R(M)\) 中由全部
\[
u\otimes v-v\otimes u,\qquad u,v\in M,
\]
生成的双边理想。商代数
\[
\operatorname{Sym}_R(M)\coloneqq T_R(M)/J
\]
称为 \(M\) 的\term[duichen-daishu]{对称代数}[symmetric algebra]。由齐次二次关系及\cref{b2:prop:homogeneous-ideal-quotient}，它有自然分次；对每个 \(d\ge0\)，其次数 \(d\) 的分量记为 \(\operatorname{Sym}_R^d(M)\)，称为 \(M\) 的 \(d\) 次\term[duichen-mi]{对称幂}[symmetric power]。张量 \(m_1\otimes\cdots\otimes m_d\) 的像记为 \(m_1\cdots m_d\)。
\end{definition}
\symidx{\(\operatorname{Sym}_R(M)\)：模的对称代数}
\symidx{\(\operatorname{Sym}_R^d(M)\)：模的 \(d\) 次对称幂}

关系在次数二，故次数零与一次部分仍分别为 \(R,M\)。全部一次生成元在商中彼此交换；任意两个字可逐个交换其中的生成元，故整个商代数交换。任何从 \(T_R(M)\) 到交换 \(R\) 代数的同态都将这些交换关系送到零，因而唯一经过这个商。所以 \(\operatorname{Sym}_R(M)\) 是强制一次生成元彼此交换得到的普遍交换商。一般模中的线性关系仍予保留，并没有因交换乘法而消失。

\begin{theorem}[对称代数的泛性质]\label{b2:thm:symmetric-algebra-universal}
设 \(R\) 为交换含幺环，\(M\) 为幺性 \(R\) 模。对任意交换含幺 \(R\) 代数 \(A\) 及 \(R\) 线性映射 \(f:M\rightarrow A\)，存在唯一保幺 \(R\) 代数同态
\[
\widehat f:\operatorname{Sym}_R(M)\longrightarrow A
\]
在一次分量上等于 \(f\)。它将 \(m_1\cdots m_d\) 送到 \(f(m_1)\cdots f(m_d)\)。
\end{theorem}
\begin{proof}
\cref{b2:thm:tensor-algebra-universal}先把 \(f\) 延拓为 \(T_R(M)\rightarrow A\)。每个关系的像为 \(f(u)f(v)-f(v)f(u)=0\)，因为 \(A\) 交换。因此延拓同态将 \(J\) 送到零，诱导唯一商同态。反过来，商中的一次分量和结构标量生成全代数，所需同态在它们上的值已指定，故唯一。
\end{proof}

对线性映射 \(f:M\rightarrow N\)，上述定理给出分次同态 \(\operatorname{Sym}_R(f)\)，其在 \(d\) 次部分的公式为
\[
\operatorname{Sym}_R^d(f)(m_1\cdots m_d)=f(m_1)\cdots f(m_d).
\]
这些映射保持恒等与复合，因为在一次生成元上保持。因此对称代数与每次对称幂都是函子。目标交换的条件不能删去：若 \(M=R^2\)，把两个基向量任意映到非交换代数中的两个元素，只有它们交换时才可经 \(\operatorname{Sym}_R(M)\) 延拓。

\ProseParagraphBreak
在次数 \(d\) 中，乘积 \(m_1\cdots m_d\) 不再记录各因子的排列。因此要把张量与它的置换像识别，也就是商去它们的差；这不同于从张量幂中选出已经被每个置换固定的元素。

\begin{proposition}\label{b2:prop:symmetric-power-universal}
设 \(R\) 为交换含幺环，\(M,N\) 为幺性 \(R\) 模，\(d\ge1\)。令 \(S_d\) 通过置换因子作用于 \(M^{\otimes d}\)，\(q_d:M^{\otimes d}\rightarrow\operatorname{Sym}_R^d(M)\) 为取商映射，则 \(q_d\) 的核由所有 \(t-P_\sigma t\)（\(t\in M^{\otimes d}\)、\(\sigma\in S_d\)）生成。因此 \(q_d\) 给出自然同构
\[
\operatorname{Sym}_R^d(M)\cong
M^{\otimes d}/\langle t-P_\sigma t\mid t\in M^{\otimes d},\ \sigma\in S_d\rangle_R.
\]
从 \(\operatorname{Sym}_R^d(M)\) 到 \(N\) 的 \(R\) 线性映射，与满足
\[
b(m_{\sigma(1)},\ldots,m_{\sigma(d)})=b(m_1,\ldots,m_d)
\quad(\sigma\in S_d,\ m_i\in M)
\]
的多线性映射 \(b:M^d\rightarrow N\) 一一对应。这样的 \(b\) 称为\term[duichen-duoxianxing-yingshe]{对称多线性映射}[symmetric multilinear map]。
\end{proposition}
\begin{proof}
设 \(K_d\) 为所述差生成的子模。先看相邻两个位置的交换：在纯张量上，\(t-P_{(i\ i+1)}t\) 是一个二次交换关系左右各乘其余纯张量的结果，因此在 \(J\) 中。任意置换均可通过相邻交换完成。若这些交换依次给出纯张量 \(t=t_0,t_1,\ldots,t_r=P_\sigma t\)，则每个 \(t_{j-1}-t_j\) 都属于 \(J\)，而 \(t-P_\sigma t=\sum_{j=1}^r(t_{j-1}-t_j)\)，故它也属于 \(J\)。这先对纯张量成立，再由线性对所有 \(t\) 成立，故 \(K_d\subseteq J\cap M^{\otimes d}\)。

反向，将 \(J\) 的元素写成 \(a(u\otimes v-v\otimes u)b\) 的有限和，再取次数 \(d\) 的分量，并把 \(a,b\) 展开成纯张量的有限和。所得每项都是某个纯张量与相邻交换后的张量之差，所以属于 \(K_d\)。因此核恰为 \(K_d\)。对任意 \(R\) 线性映射 \(f:M\to M'\)，有 \(f^{\otimes d}P_\sigma=P_\sigma f^{\otimes d}\)，故该商同构与 \(\operatorname{Sym}_R^d(f)\) 相容，确为自然同构。

多线性映射先由\cref{b2:thm:multiple-tensor-universal}唯一合成 \(M^{\otimes d}\rightarrow N\)。变量置换不改变值，恰好要求它在上述核上为零，因而唯一经 \(q_d\) 分解。反向复合显然逐变量线性且对称。
\end{proof}

\subsection{有限自由模的多项式解释}
\label{b2:subsec:0803:02}

\begin{theorem}[多项式解释]\label{b2:thm:symmetric-polynomial-algebra}
设 \(R\) 为交换含幺环，\(M\) 为以 \(e_1,\ldots,e_r\) 为基的幺性有限自由 \(R\) 模，则存在唯一分次 \(R\) 代数同构
\[
\operatorname{Sym}_R(M)\longrightarrow R[X_1,\ldots,X_r],
\qquad e_i\longmapsto X_i.
\]
\end{theorem}
\begin{proof}
将基向量送到不定元给出线性映射 \(M\rightarrow R[X_1,\ldots,X_r]\)。目标交换，故\cref{b2:thm:symmetric-algebra-universal}给出正向同态。反向把多项式 \(\sum_\alpha a_\alpha X^\alpha\) 送到 \(\sum_\alpha a_\alpha e_1^{\alpha_1}\cdots e_r^{\alpha_r}\)。和为有限和，商中的 \(e_i\) 彼此交换，因此多项式的加法与乘法公式保证它是代数同态。两复合分别在各 \(e_i\) 与 \(X_i\) 上为恒等，并在结构标量上为恒等；生成性使两复合在全代数为恒等。两映射保持生成元的次数，故保持分次。
\end{proof}

\begin{proposition}\label{b2:prop:symmetric-power-basis}
设 \(R\) 为交换含幺环，\(M\) 为以 \(e_1,\ldots,e_r\) 为基的幺性有限自由 \(R\) 模，\(d\ge0\)，则次数 \(d\) 的单项式
\[
e_1^{a_1}\cdots e_r^{a_r},\qquad a_i\ge0,\quad a_1+\cdots+a_r=d,
\]
构成 \(\operatorname{Sym}_R^d(M)\) 的一组基。若 \(R\ne0\)、\(r\ge1\)，则其秩为 \(\binom{r+d-1}{d}\)。当 \(r=0\) 时，零次部分为 \(R\)，所有正次部分为零。
\end{proposition}
\begin{proof}
多项式作为有限形式和，其不同单项式的系数唯一；上一同构把全次数 \(d\) 的单项式送到题述元素，所以它们构成基。计算数目时，把 \(d\) 个相同的记号排成一行，并插入 \(r-1\) 道隔板。第 \(i\) 段记号的个数就是 \(a_i\)，允许空段。这与从 \(r+d-1\) 个位置选出 \(d\) 个记号位置一一对应，故数目为所述二项式系数。零模情形直接由定义得到。
\end{proof}

例如 \(R^2\) 的三次对称幂以 \(e_1^3,e_1^2e_2,e_1e_2^2,e_2^3\) 为基，秩为四；相应三次张量幂则有八个有序纯基张量。对称商把 \(e_1\otimes e_1\otimes e_2\) 的三个不同排列都映到 \(e_1^2e_2\)；只有把这三个张量先相加，其像才是 \(3e_1^2e_2\)。单个张量的像不需要乘上排列数。

\ProseParagraphBreak
对有限维 \(k\) 向量空间 \(V\)，\(\operatorname{Sym}_k(V)\) 的一次生成元是向量，而 \(V\) 上的线性坐标函数属于对偶空间 \(V^*\)。因此，由线性坐标函数生成的形式多项式代数应写为 \(\operatorname{Sym}_k(V^*)\)：线性泛函 \(\lambda\) 在 \(v\) 处取值为 \(\lambda(v)\)，乘积 \(\lambda_1\cdots\lambda_d\) 则取值为 \(\prod_j\lambda_j(v)\)。选定 \(V\) 的基后，其对偶基对应熟悉的坐标不定元。两空间虽同维，\(V\) 与 \(V^*\) 却没有由向量空间本身指定的识别，不能把向量生成元直接当作坐标函数。

\ProseParagraphBreak
形式多项式也不等同于它定义的函数。例如在 \(\mathbb F_p\) 上，\(X^p-X\) 是非零多项式，却在每个域元素处取零。对称代数采用形式多项式的乘法和系数，不因所有取值为零就把它的元素置零；\secref{b2:sec:0804}将说明这一区别如何影响不变环的定义。

\subsection{对称张量、对称商与对称化}
\label{b2:subsec:0803:03}

先考虑一个交换两个基向量的作用：在自由模 \(Ru\oplus Rv\) 中，不变子模为 \(R(u+v)\)，余不变量则为 \((Ru\oplus Rv)/R(u-v)\)。前者要求两个坐标本来相同，后者让 \(u,v\) 的类相同。下面把这两种构造用于张量因子的置换。

\ProseParagraphBreak
设 \(R\) 为交换含幺环，\(M\) 为幺性 \(R\) 模，\(d\ge1\)。令 \(S_d\) 按\cref{b2:prop:multiple-tensor-functoriality}的公式作用在 \(M^{\otimes d}\) 上。这一作用的\term[bubianliang]{不变量}[invariants]是子模
\[
(M^{\otimes d})^{S_d}
=\{\,t\in M^{\otimes d}\mid P_\sigma t=t\;\text{对所有}\;\sigma\in S_d\,\},
\]
其中的元素称为\term[duichen-zhangliang]{对称张量}[symmetric tensor]。同一作用的\term[yububianliang]{余不变量}[coinvariants]是商模
\[
(M^{\otimes d})_{S_d}
\coloneqq M^{\otimes d}/\langle t-P_\sigma t\mid t\in M^{\otimes d},\ \sigma\in S_d\rangle_R.
\]
由\cref{b2:prop:symmetric-power-universal}，商映射给出 \(\operatorname{Sym}_R^d(M)\cong(M^{\otimes d})_{S_d}\)。因此包含与取商组成
\[
(M^{\otimes d})^{S_d}\hookrightarrow M^{\otimes d}
\xrightarrow{q_d}\operatorname{Sym}_R^d(M)
\cong(M^{\otimes d})_{S_d},
\]
其中 \(q_d\) 满射，总有指定映射
\[
q_d\big|_{(M^{\otimes d})^{S_d}}:
(M^{\otimes d})^{S_d}\longrightarrow\operatorname{Sym}_R^d(M).
\]
线性映射的张量幂保持不变子模，并与 \(q_d\) 相容，所以这个限制映射是自然的。问题在于它何时成为同构，而不是两个模是否偶然具有相同的秩或维数。
\symidx{\((M^{\otimes d})^{S_d},(M^{\otimes d})_{S_d}\)：对称张量的不变量与余不变量}

\ProseParagraphBreak
在次数二，若 \(2\) 可逆，令 \(P\) 为交换两个因子的算子，则 \((t+Pt)/2\) 被 \(P\) 固定；若 \(t\) 本来对称，这个平均仍为 \(t\)。一般次数下，对全部 \(d!\) 个置换同样取平均，就能得到一个固定不变张量的投影。

\begin{proposition}\label{b2:prop:symmetrization-characteristic}
设 \(R\) 为交换含幺环，\(M\) 为幺性 \(R\) 模，\(d\ge1\)。若 \(d!\) 在 \(R\) 中可逆，则\term[duichen-hua]{对称化}[symmetrization]算子
\[
E_d=\frac1{d!}\sum_{\sigma\in S_d}P_\sigma
\]
是投影到 \((M^{\otimes d})^{S_d}\) 的线性映射，并诱导自然同构
\[
\overline E_d:\operatorname{Sym}_R^d(M)\xrightarrow{\ \cong\ }(M^{\otimes d})^{S_d}.
\]
其逆是商映射 \(q_d\) 在不变子模上的限制。
\end{proposition}
\begin{proof}
左乘一个固定置换只重新排列求和项，因此 \(P_\pi E_d=E_d\)，像都不变；若 \(t\) 已不变，则 \(E_dt=t\)。所以像正是不变子模，且 \(E_d^2=E_d\)。

同样，\(E_dP_\pi=E_d\)，故 \(E_d\) 零化全部 \(t-P_\pi t\)。由\cref{b2:prop:symmetric-power-universal}，它唯一经 \(q_d\) 分解为 \(\overline E_d\)。在对称商中 \(q_dP_\sigma=q_d\)，因而 \(q_dE_d=q_d\)，说明 \(q_d\overline E_d\) 为恒等；在不变子模上 \(E_d\) 为恒等，又说明 \(\overline E_dq_d\) 为恒等。线性映射的张量幂与每个 \(P_\sigma\) 交换，所以也与平均算子相容；所得同构因此自然。
\end{proof}

阶乘可逆保证上述同构，但它不可逆时须另行检查，不可据此一概断言同构失败。特征零也不保证阶乘可逆：取 \(R=\mathbb Z\)、\(M=\mathbb Ze\oplus\mathbb Zf\)。交换两个因子的不变张量恰为
\[
a(e\otimes e)+b(e\otimes f+f\otimes e)+c(f\otimes f),
\qquad a,b,c\in\mathbb Z,
\]
因为在纯基张量的展开中，\(e\otimes f\) 与 \(f\otimes e\) 的系数必须相同。它在 \(\operatorname{Sym}_{\mathbb Z}^2(M)\) 中的像为
\[
ae^2+2bef+cf^2.
\]
\(e^2,ef,f^2\) 是这个对称幂的一组基，所以 \(ef\) 不在像中，否则需要整数 \(b\) 满足 \(2b=1\)。这里商映射在不变子模上的限制不是满射，障碍是 \(2\) 在 \(\mathbb Z\) 中不可逆。

\ProseParagraphBreak
正特征下还可能出现单射性失败。取 \(k\) 特征为素数 \(p\)，\(V=ke\oplus kf\)，在 \(V^{\otimes p}\) 中令
\[
s=\sum_{j=1}^p e\otimes\cdots\otimes
\underset{\text{第 }j\text{ 个位置}}{f}
\otimes\cdots\otimes e.
\]
各项是互不相同的纯基张量，所以 \(s\ne0\)；置换只重排这些项，所以 \(s\) 不变。但
\[
q_p(s)=p\,e^{p-1}f=0.
\]
因此商映射在不变张量上不是单射。最小例子是特征二下的 \(e\otimes f+f\otimes e\)：它非零且对称，却在二次对称幂中为零。反过来，\(ef\) 仍是对称幂的一个非零基元素。

\begin{example}\label{b2:exm:cubic-symmetrization-orbits}
设 \(k\) 是域，\(V=ke\oplus kf\)。在 \(V^{\otimes3}\) 的八个纯基张量上，\(S_3\) 的置换作用有四个轨道，按 \(f\) 出现零、一、二、三次分类。相应轨道和为
\[
\begin{aligned}
s_0&=e\otimes e\otimes e,\\
s_1&=e\otimes e\otimes f+e\otimes f\otimes e+f\otimes e\otimes e,\\
s_2&=e\otimes f\otimes f+f\otimes e\otimes f+f\otimes f\otimes e,\\
s_3&=f\otimes f\otimes f.
\end{aligned}
\]
置换作用只重排纯基张量，因而一个张量不变，当且仅当同一轨道上的系数相同。四个轨道的支集互不相交，所以 \(s_0,s_1,s_2,s_3\) 线性无关，并生成全部不变张量。它们是一组基；\(\operatorname{Sym}_k^3(V)\) 则以 \(e^3,e^2f,ef^2,f^3\) 为基。限制商映射把 \(s_0,s_1,s_2,s_3\) 分别送到 \(e^3,3e^2f,3ef^2,f^3\)，因此相对于这两组基的矩阵为
\[
\operatorname{diag}(1,3,3,1).
\]
若 \(k=\mathbb Q\)，它是同构，且 \(e^2f\) 的逆像为 \(s_1/3\)。这也等于六个置换的平均：\(e\otimes e\otimes f\) 的每个不同排列出现两次，所以平均为 \(2s_1/6=s_1/3\)。

若 \(\operatorname{char}k=2\)，上述矩阵为单位矩阵，限制映射仍是同构，尽管 \(3!=0\) 不可逆。若 \(\operatorname{char}k=3\)，它的核为 \(ks_1\oplus ks_2\)，像为 \(ke^3\oplus kf^3\)，秩为二；两端都是四维，也不保证指定映射可逆。因此阶乘可逆是统一保证同构的充分条件，并不是一个特定模上限制映射可逆的必要条件。
\end{example}

删去平均公式中的分母，就得到始终有定义的求和算子 \(N_d=\sum_\sigma P_\sigma\)。但它可能把本来的不变张量也送到零。例如在特征 \(p\) 的域 \(k\) 上，一维空间 \(ke\) 的次数 \(p\) 张量幂中，全部置换作用都是恒等，故 \(N_p=p!\operatorname{id}=0\)，而张量幂、不变子空间和对称幂都非零。因此只删去分母不能补救平均公式；这里限制商映射反而是同构，失败的是这个求和算子。

\ProseParagraphBreak
正特征不仅影响平均过程，还使对称代数内部的取幂运算出现新的加法性。设 \(p\) 为素数，\(R\) 为特征 \(p\) 的交换含幺环，\(M\) 为幺性 \(R\) 模。对称代数交换，故
\[
(v+w)^p=v^p+w^p,\qquad (av)^p=a^pv^p.
\]
第一式由二项式系数 \(\binom pi\) 对 \(0<i<p\) 都被 \(p\) 整除得到。因此 \(\Phi_p:M\rightarrow\operatorname{Sym}_R^p(M)\)、\(v\mapsto v^p\) 是可加映射，并对 Frobenius 环同态 \(F_R:R\rightarrow R\)、\(a\mapsto a^p\) 半线性：\(\Phi_p(av)=F_R(a)\Phi_p(v)\)。这里半线性指映射可加，而标量 \(a\) 须先经过 \(F_R\) 才作用于像；若 \(F_R\) 为恒等，就得到通常的 \(R\) 线性。
\symidx{\(F_R\)：特征 \(p\) 环上的 Frobenius 同态}

\ProseParagraphBreak
若进一步 \(R=k\) 为域且 \(M\ne0\)，对非零 \(v\in M\)，先在 \(kv\) 上定义 \(av\mapsto a\)，再由\cref{b2:lem:functional-extension}延拓为 \(\ell:M\to k\)，便有 \(\ell(v)=1\)。对称多线性映射 \((m_1,\ldots,m_p)\mapsto\prod_i\ell(m_i)\) 由\cref{b2:prop:symmetric-power-universal}诱导线性映射 \(\operatorname{Sym}_k^p(M)\to k\)，将 \(v^p\) 送到 \(1\)，故 \(v^p\ne0\)。若 \(\Phi_p\) 是 \(k\) 线性的，便有 \((a^p-a)v^p=0\) 对每个 \(a\in k\) 成立；域上的非零向量不能被非零标量零化，故 \(a^p=a\)。因此在这个域且模非零的范围内，\(\Phi_p\) 线性当且仅当 \(F_k\) 为恒等。一般环上不能由 \(\Phi_p\) 线性就这样消去 \(v^p\)。因子对称、取对称商和取 \(p\) 次幂是相互联系但不同的运算。

\ProseParagraphBreak
第六章把扩张标量写成 \(S\otimes_R-\)。现在的代数构造只用有限张量、直和与生成关系，改变底环后也能与扩张标量相容；不必为模选择一组基。

\begin{proposition}\label{b2:prop:tensor-symmetric-base-change}
设 \(R,S\) 是交换含幺环，\(\varphi:R\to S\) 是保幺环同态，\(M\) 是幺性 \(R\) 模，记 \(P=S\otimes_RM\)。存在自然的分次 \(S\) 代数同构
\[
\begin{aligned}
\alpha_T&:S\otimes_RT_R(M)\xrightarrow{\ \cong\ }T_S(P),\\
\alpha_{\mathrm{Sym}}&:S\otimes_R\operatorname{Sym}_R(M)
\xrightarrow{\ \cong\ }\operatorname{Sym}_S(P).
\end{aligned}
\]
左侧的乘法均由 \((s\otimes a)(t\otimes b)=st\otimes ab\) 给出，次数 \(d\) 的分量分别为 \(S\otimes_RM^{\otimes_Rd}\) 与 \(S\otimes_R\operatorname{Sym}_R^d(M)\)。零次的同构均为 \(S\otimes_RR\to S\)、\(s\otimes r\mapsto s\varphi(r)\)。对 \(d\ge1\)，其公式为
\[
\begin{aligned}
\alpha_T\bigl(s\otimes(m_1\otimes_R\cdots\otimes_Rm_d)\bigr)
&=(s\otimes_Rm_1)\otimes_S(1\otimes_Rm_2)
\otimes_S\cdots\otimes_S(1\otimes_Rm_d),\\
\alpha_{\mathrm{Sym}}\bigl(s\otimes(m_1\cdots m_d)\bigr)
&=(s\otimes_Rm_1)(1\otimes_Rm_2)\cdots(1\otimes_Rm_d).
\end{aligned}
\]
当 \(d=1\) 时，右侧只保留首因子。不要求 \(S\) 作为 \(R\) 模平坦。
\end{proposition}
\begin{proof}
先给左侧代数结构。对任意含幺 \(R\) 代数 \(A\)，因 \(R\) 的结构像在 \(A\) 中心，所述乘法保持两处 \(R\) 平衡关系。例如
\[
(s\varphi(r)\otimes a)(t\otimes b)
=st\varphi(r)\otimes ab
=st\otimes (ra)b
=(s\otimes ra)(t\otimes b),
\]
另一处由 \(arb=abr\) 同样核对。因此乘法良定义；结合律与单位元 \(1\otimes1\) 由 \(S,A\) 的对应等式继承，\(s\mapsto s\otimes1\) 的像在中心，故得到 \(S\) 代数。若 \(A\) 分次，\cref{b2:prop:tensor-direct-sums}给出 \(S\otimes_RA\cong\bigoplus_d(S\otimes_RA_d)\)，且上述乘法使次数相加。

令 \(\iota_R:M\to T_R(M)\)、\(\iota_S:P\to T_S(P)\) 为一次分量包含。把 \(T_S(P)\) 经 \(\varphi\) 看成 \(R\) 代数，映射 \(m\mapsto\iota_S(1\otimes m)\) 为 \(R\) 线性：平衡关系 \(1\otimes rm=\varphi(r)\otimes m\) 保证标量相容。\cref{b2:thm:tensor-algebra-universal}给出 \(R\) 代数同态 \(h:T_R(M)\to T_S(P)\)。公式
\[
\alpha_T(s\otimes a)=s\cdot h(a)
\]
对 \(R\) 平衡，因为 \(h(ra)=\varphi(r)h(a)\)，所以诱导 \(S\) 线性映射。\(S\) 标量在目标中心，故
\[
\alpha_T\bigl((s\otimes a)(t\otimes b)\bigr)
=st\cdot h(ab)
=(s\cdot h(a))(t\cdot h(b)),
\]
并且它保单位元。\(h\) 在一次生成元上的值给出题述纯张量公式；将标量 \(s\) 移入第一个 \(P\) 因子即得到所写形式。

反向，\(S\) 线性映射
\[
j:P\to S\otimes_RT_R(M),\qquad s\otimes m\longmapsto s\otimes\iota_R(m)
\]
由 \(\operatorname{id}_S\otimes\iota_R\) 给出，故良定义。再次用张量代数的泛性质，得到 \(S\) 代数同态 \(\beta_T:T_S(P)\to S\otimes_RT_R(M)\)。它在纯张量上的公式为
\[
\beta_T\bigl((s_1\otimes m_1)\otimes_S\cdots\otimes_S(s_d\otimes m_d)\bigr)
=(s_1\cdots s_d)\otimes(m_1\otimes_R\cdots\otimes_Rm_d),
\]
在零次部分则为 \(s\mapsto s\otimes1\)。两个复合都固定 \(S\) 标量；\(\alpha_T\beta_T\) 在每个一次生成元 \(s\otimes m\in P\) 上为恒等，\(\beta_T\alpha_T\) 在 \(s\otimes\iota_R(m)\) 上为恒等。这些元素连同 \(S\) 生成相应代数，故两复合在全代数上为恒等。两向映射把一次生成元送到一次分量，把标量送到零次分量，所以保持分次。

对称情形中，把 \(m\) 送到 \(\operatorname{Sym}_S(P)\) 的一次生成元 \(1\otimes m\)，由\cref{b2:thm:symmetric-algebra-universal}延拓为 \(R\) 代数同态 \(h_{\mathrm{Sym}}:\operatorname{Sym}_R(M)\to\operatorname{Sym}_S(P)\)。再定义 \(\alpha_{\mathrm{Sym}}(s\otimes a)=s\cdot h_{\mathrm{Sym}}(a)\)；与张量情形相同的平衡等式和乘法等式说明它是 \(S\) 代数同态，且给出题述公式。

反向，\(S\otimes_R\operatorname{Sym}_R(M)\) 是交换 \(S\) 代数。将 \(s\otimes m\in P\) 送到目标代数一次分量 \(S\otimes_R\operatorname{Sym}_R^1(M)\) 中的对应元素，得到 \(S\) 线性映射；再用对称代数的泛性质，得到 \(\beta_{\mathrm{Sym}}:\operatorname{Sym}_S(P)\to S\otimes_R\operatorname{Sym}_R(M)\)。具体地，
\[
\beta_{\mathrm{Sym}}\bigl((s_1\otimes m_1)\cdots(s_d\otimes m_d)\bigr)
=(s_1\cdots s_d)\otimes(m_1\cdots m_d).
\]
两个复合在 \(S\) 标量和一次生成元上都为恒等，故在全代数上为恒等；生成元的次数也得到保留，因而同构分次。这里所用的是两个交换代数的泛性质，并未要求张量积保持单射。

对任意 \(R\) 线性映射 \(f:M\to M'\)，先扩张标量再作用于一次生成元，与先作用于 \(m\) 再按上述公式映射，都得到 \(s\otimes f(m)\)；在 \(S\) 标量上也相同。生成性因此给出
\[
\begin{aligned}
\alpha_{T,M'}(\operatorname{id}_S\otimes T_R(f))
&=T_S(\operatorname{id}_S\otimes f)\alpha_{T,M},\\
\alpha_{\mathrm{Sym},M'}(\operatorname{id}_S\otimes\operatorname{Sym}_R(f))
&=\operatorname{Sym}_S(\operatorname{id}_S\otimes f)\alpha_{\mathrm{Sym},M}.
\end{aligned}
\]
这证明两个同构对 \(M\) 自然。若 \(S'\) 也是交换含幺 \(R\) 代数，并给定保幺 \(R\) 代数同态 \(\psi:S\to S'\)，两侧由 \(s\mapsto\psi(s)\)、\(s\otimes m\mapsto\psi(s)\otimes m\) 诱导的映射，也与这些同构相容：上述纯张量公式两种复合均只将系数 \(s\) 换成 \(\psi(s)\)，生成性使相容性在全代数上成立。
\end{proof}

% !TeX root = ../../Book2.tex
% 纲要标题：### 8.4 张量运算与表示
% 本轮补充的纲要细目：
% * 由对偶空间的对称代数定义形式多项式上的群作用，写清逆向作用、不变环及有限域上与点集函数的区别
% * 以反射及单项式权计算不变环，不将逐次不动空间的计算自动视为全部生成关系已知

\section{张量运算与表示}
\label{b2:sec:0804}

若群同时作用于几个向量空间，多线性构造也随之得到群作用。这个过程使一个表示产生新的表示，并把同态问题转为不变张量问题。本节以任意域 \(k\) 和任意群 \(G\) 开始；凡需有限维或特征零处都另行说明。因子置换由 \(S_d\) 作用，表示的不变量由 \(G\) 作用；这两个群作用应始终区分。

% 纲要要点（供目录核对）：
%
% * 对偶、张量积与对称幂表示
% * 不变量和对称张量的区别
% * 第五卷 张量表示与 Schur–Weyl 理论
%
% ---
%

\subsection{对偶、张量积与对称幂表示}
\label{b2:subsec:0804:01}

\begin{definition}\label{b2:def:linear-representation}
设 \(k\) 为域，\(G\) 为群，\(V\) 为 \(k\) 向量空间。如果给定群同态
\[
\rho_V:G\longrightarrow\operatorname{GL}_k(V),
\]
就称 \(\rho_V\) 为 \(G\) 在 \(V\) 上的\term[xianxing-biaoshi]{线性表示}[linear representation]。
常把 \(\rho_V(g)v\) 写成 \(gv\)。在本节给定的表示中，显式作用记号 \(g\cdot v\) 也表示同一个 \(\rho_V(g)v\)；对另一个表示 \(W\)，则按它自己的 \(\rho_W\) 解释。若 \(V,W\) 均为 \(G\) 的 \(k\) 线性表示，\(k\) 线性映射 \(f:V\rightarrow W\) 满足 \(f(gv)=g\cdot f(v)\) 对每个 \(g\in G\)、\(v\in V\) 成立，则称 \(f\) 为表示同态或等变线性映射；全部这类映射组成空间 \(\operatorname{Hom}_G(V,W)\)。
\end{definition}
\symidx{\(\operatorname{Hom}_G(V,W)\)：表示之间的等变线性映射空间}
\symidx{\(\rho_V(g)\)：群元素在表示空间上的线性作用}

群作用公理在这里具体为 \(1v=v\)、\((gh)v=g(hv)\)，并要求每个 \(g\) 的作用可逆且线性。定义没有要求表示忠实，也没有要求群有限。例如无穷循环群由一个生成元 \(g\) 生成，在 \(\mathbb Q^2\) 上可以规定 \(g\) 的作用矩阵为 \(\operatorname{diag}(2,1/2)\)，再由整数次幂定义全部群元素的作用。

\ProseParagraphBreak
向量与线性泛函同时变换时，自然的求值配对应保持不变，也就是要求 \((g\lambda)(gv)=\lambda(v)\)，其中 \(k\) 上采用平凡作用。把 \(gv\) 记为 \(w\)，这个要求便强迫 \((g\lambda)(w)=\lambda(g^{-1}w)\)。因此对偶作用中的逆元来自求值配对的相容性：向量的变换与泛函的变换必须互相抵消。

\begin{proposition}\label{b2:prop:tensor-dual-symmetric-representations}
设 \(k\) 为域，\(G\) 为群，\(V,W\) 为 \(G\) 的 \(k\) 线性表示，\(d\ge0\)。以下公式分别定义 \(V^*\)、\(V\otimes_kW\) 及 \(\operatorname{Sym}_k^d(V)\) 上的表示：
\[
\begin{aligned}
(g\lambda)(v)&=\lambda(g^{-1}v),\\
g(v\otimes w)&=gv\otimes gw,\\
g(v_1\cdots v_d)&=(gv_1)\cdots(gv_d).
\end{aligned}
\]
它们分别称为\term[duiou-biaoshi]{对偶表示}[dual representation]、\term[zhangliangji-biaoshi]{张量积表示}[tensor product representation]和\term[duichenmi-biaoshi]{对称幂表示}[symmetric power representation]。
\end{proposition}
\begin{proof}
对偶公式对 \(v\) 线性，也对 \(\lambda\) 线性。对 \(g,h\in G\)，
\[
\bigl(g(h\lambda)\bigr)(v)
=(h\lambda)(g^{-1}v)
=\lambda(h^{-1}g^{-1}v)
=\bigl((gh)\lambda\bigr)(v).
\]
单位元作用为恒等，\(g^{-1}\) 的作用为逆，故为表示。使用逆元正是为了抵消对偶映射反转复合次序的效应。

张量公式由 \(\rho_V(g)\otimes\rho_W(g)\) 给出，良定义来自\cref{b2:prop:tensor-maps}；该命题的复合公式保证群乘法，逆由 \(g^{-1}\) 给出。对称幂的映射是 \(\operatorname{Sym}_k^d\bigl(\rho_V(g)\bigr)\)，对称幂函子保持恒等和复合，因此同样为表示。也可直接观察：张量幂作用将每个交换关系送到另一个交换关系，所以它通过对称商。
\end{proof}

对偶构造把特征值取逆，张量构造则把各因子的特征值相乘。具体地，若 \(g\) 在有限维 \(V\) 的基 \(e_1,\ldots,e_r\) 上对角作用，\(ge_i=a_i e_i\)，则它在对偶基上的特征值为 \(a_i^{-1}\)，在张量基 \(e_{i_1}\otimes\cdots\otimes e_{i_d}\) 上的特征值为 \(a_{i_1}\cdots a_{i_d}\)，在对称基 \(e_1^{b_1}\cdots e_r^{b_r}\) 上的特征值为 \(a_1^{b_1}\cdots a_r^{b_r}\)。这里各 \(a_i\ne0\)，因为群元素作用可逆。若一个张量被 \(g\) 固定，其基展开中只能出现上述特征值等于一的项。

\ProseParagraphBreak
\cref{b2:prop:finite-free-hom-tensor}已在有限自由模上建立求值同构；有限维向量空间当然满足该条件。给 \(V,W\) 加上群作用后，还要使这个同构与作用相容。对线性映射 \(A:V\to W\)，变换输入时先由 \(g^{-1}\) 拉回，施加 \(A\) 后再由 \(g\) 推进输出，这给出下面的 Hom 作用。

\begin{proposition}\label{b2:prop:hom-as-tensor-representation}
设 \(k\) 为域，\(V\) 为有限维 \(k\) 向量空间，\(W\) 为任意 \(k\) 向量空间，则
\[
\Theta:V^*\otimes_kW\longrightarrow\operatorname{Hom}_k(V,W),
\qquad
\lambda\otimes w\longmapsto\bigl(v\mapsto\lambda(v)w\bigr)
\]
是自然线性同构。若群 \(G\) 在 \(V,W\) 上均有 \(k\) 线性表示，则它是 \(G\) 表示同构，其中右端的作用为
\[
(g\cdot A)(v)=\rho_W(g)\bigl(A(\rho_V(g)^{-1}v)\bigr).
\]
\end{proposition}
\begin{proof}
取\cref{b2:prop:finite-free-hom-tensor}中的 \(R=k\)、\(P=V\)、\(M=W\)，即得所写自然线性同构。该命题还给出坐标逆式
\[
\Theta^{-1}(A)=\sum_i e_i^*\otimes A(e_i),
\]
其中 \(e_i\) 是 \(V\) 的任意基，\(e_i^*\) 是其对偶基；该表达式所确定的张量不依赖所选基。

群作用的相容性直接在纯张量上核对：
\[
\Theta(g\lambda\otimes gw)(v)
=\lambda(\rho_V(g)^{-1}v)\,\rho_W(g)w
=\rho_W(g)\bigl(\Theta(\lambda\otimes w)(\rho_V(g)^{-1}v)\bigr).
\]
这就是右端所给的作用。它满足群乘法公理，因为对 \(A:V\to W\) 先作用 \(h\) 再作用 \(g\)，所得为 \(\rho_W(gh)A\rho_V(gh)^{-1}\)。
\end{proof}

当 \(W=V\) 时，恒等算子对应的张量 \(\sum_i e_i^*\otimes e_i\) 不依赖基，并被所有 \(g\) 固定，因为 \(\rho_V(g)\operatorname{id}_V\rho_V(g)^{-1}=\operatorname{id}_V\)。向量与对偶向量的变换互相抵消，并不要求每个基向量被固定。一般算子的矩阵也可按列作这种转换：若 \(V\) 有基 \(e_1,e_2\)，且 \(A\) 的矩阵为 \(\begin{pmatrix}a&b\\c&d\end{pmatrix}\)，则
\[
\Theta^{-1}(A)
=e_1^*\otimes(ae_1+ce_2)+e_2^*\otimes(be_1+de_2).
\]
恒等矩阵便给出 \(e_1^*\otimes e_1+e_2^*\otimes e_2\)。

\subsection{不变量与对称张量}
\label{b2:subsec:0804:02}

\begin{definition}\label{b2:def:invariant-vector}
设 \(k\) 为域，\(G\) 为群，\(V\) 为 \(G\) 的 \(k\) 线性表示。若 \(v\in V\) 满足 \(gv=v\) 对每个 \(g\in G\) 成立，就称 \(v\) 为 \(G\) 的\term[bubian-xiangliang]{不变向量}[invariant vector]。全部不变向量组成的子空间记为
\[
V^G=\{\,v\in V\mid gv=v\;\text{对每个}\;g\in G\,\}.
\]
若子空间 \(U\subseteq V\) 满足 \(gU\subseteq U\) 对每个 \(g\in G\) 成立，就称 \(U\) 为 \(G\) 不变子空间，也称稳定子空间；这只要求作用保留整个子空间，不能因此断言其中每个向量都被逐点固定。
\end{definition}
\symidx{\(V^G\)：群作用下逐点不变的向量空间}

不变条件由一族线性方程给出，故 \(V^G\) 确为子空间。对任意两个 \(G\) 表示 \(V,W\)，上一命题所写的公式都在 \(\operatorname{Hom}_k(V,W)\) 上定义作用，不要求 \(V\) 有限维。由此有
\[
\operatorname{Hom}_k(V,W)^G=\operatorname{Hom}_G(V,W).
\]
确切地说，\(A:V\to W\) 在上述作用下不变，当且仅当对每个 \(g\in G\) 都有
\[
\rho_W(g)A\rho_V(g)^{-1}=A,
\qquad\text{即}\qquad
\rho_W(g)A=A\rho_V(g).
\]
后一个等式正是等变性。因此当 \(V\) 有限维时，\cref{b2:prop:hom-as-tensor-representation}给出
\[
(V^*\otimes_kW)^G\cong\operatorname{Hom}_G(V,W).
\]
寻找等变映射，可以转为求张量表示中的固定向量。

\ProseParagraphBreak
对称张量采用的是 \(S_d\) 对因子位置的作用，它只重新排列向量；\(G\) 不变张量采用的则是 \(G\) 的对角作用，它保持位置而同时改变各位置上的向量。即使两种作用发生在同一个空间，也不是同一个条件。取 \(G=\langle g\rangle\) 为无穷循环群，\(V=\mathbb Qe\oplus\mathbb Qf\)，规定 \(ge=2e\)、\(gf=\tfrac12f\)。在 \(V^{\otimes2}\) 中，\(e\otimes e\) 对交换两因子保持不变，却被 \(g\) 乘四，因而不是 \(G\) 不变量；\(e\otimes f\) 被 \(G\) 固定，却被交换因子送到不同的基张量 \(f\otimes e\)，因而不是对称张量。

\ProseParagraphBreak
本例中
\[
(V^{\otimes2})^G=\operatorname{span}_{\mathbb Q}\{e\otimes f,f\otimes e\},
\]
因为四个纯基张量的特征值依次为 \(4,1,1,1/4\)。其中的对称部分是一维空间
\(\mathbb Q(e\otimes f+f\otimes e)\)。二次对称幂的 \(G\) 不变部分则为 \(\mathbb Q ef\)，两者由平均同构对应；该对应确实使用了二的可逆性。

\ProseParagraphBreak
固定向量也可以在对称代数的各个齐次分量中寻找。\secref{b2:sec:0803}已说明，\(V^*\) 的元素是 \(V\) 上的线性坐标函数，而 \(\operatorname{Sym}_k(V^*)\) 用交换乘法组合这些函数。把各次不变量连同乘法一起考察，便得到多项式不变环。

\begin{definition}\label{b2:def:polynomial-invariant-ring}
设 \(k\) 为域，\(V\) 为有限维 \(k\) 向量空间。称对称代数
\[
k[V]\coloneqq\operatorname{Sym}_k(V^*)
\]
为 \(V\) 上的多项式代数。若群 \(G\) 线性作用于 \(V\)，就在生成元 \(\lambda\in V^*\) 上规定 \(g\lambda=\lambda\circ g^{-1}\)，并按代数同态延拓，得到 \(G\) 在 \(k[V]\) 上的作用。被所有群元素固定的多项式称为\term[bubian-duoxiangshi]{不变多项式}[invariant polynomial]；它们组成的子代数
\[
k[V]^G=\{\,f\in k[V]\mid gf=f\;\text{对每个}\;g\in G\,\}
\]
称为这一作用的\term[bubian-huan]{不变环}[invariant ring]。
\end{definition}
\symidx{\(k[V]^G\)：表示空间上的多项式不变量环}
\symidx{\(k[V]=\operatorname{Sym}_k(V^*)\)：向量空间上的多项式代数}

\cref{b2:prop:tensor-dual-symmetric-representations}保证这个延拓满足群乘法，并保持每个齐次分量。在一组基及其对偶坐标 \(x_1,\ldots,x_n\) 下，\cref{b2:thm:symmetric-polynomial-algebra}把 \(k[V]\) 识别为 \(k[x_1,\ldots,x_n]\)。把多项式在 \(v\in V\) 处求值，便有
\[
(gf)(v)=f(g^{-1}v).
\]
这是函数沿 \(g^{-1}:V\to V\) 的拉回作用。使用逆元使拉回后的复合符合群乘法的顺序；在形式多项式代数中，它就是由线性坐标函数的对偶作用诱导的代数自同构。因而不变多项式的和与积仍不变，单位元也不变，定义中的集合确为子代数。

\ProseParagraphBreak
逐次求不动空间，只得到不变环的各个向量空间分量；要说明整个环如何生成，还须研究这些分量之间的乘法。例如，即使已知二次不变量恰为 \(kq\)，其中 \(q\ne0\)，也不能据此断言 \(k[V]^G=k[q]\)：还须知道其他次数是否出现不能由 \(q\) 的幂表示的不变量。找到多个生成元后，还要确定它们之间的关系。第五卷第22章将把这些问题与 Hilbert 级数、生成关系及商联系起来。\footnote{希尔伯特（David Hilbert，1862-1943），德国数学家，证明多项式环的基定理，并推进不变量论与几何公理化。}

\ProseParagraphBreak
在无限域上，不同多项式给出不同的点集函数。一元非零多项式的根数不超过次数；多元情形将多项式按最后一个变量展开，固定其余变量后应用一元结论，再对系数多项式归纳，即知在 \(k^n\) 上处处为零的多项式只能为零。在有限域上则不然。例如 \(k=\mathbb F_q\)、\(V=k\) 时，\(x^q-x\) 的最高次系数为一，因而是非零多项式；但 \(k^\times\) 是阶为 \(q-1\) 的群，Lagrange 定理给出每个非零 \(a\in k\) 都满足 \(a^{q-1}=1\)，加上 \(a=0\) 便得所有 \(a\) 都满足 \(a^q-a=0\)。所以这个非零多项式定义零函数。这里的 \(k[V]\) 始终指形式多项式代数；不变性要求代数中的等式，而不能一概只检查有限个点上的取值。

\newpage
\begin{example}\label{b2:exm:reflection-polynomial-invariants}
设 \(\operatorname{char}k\ne2\)，群 \(C_2=\{1,s\}\) 在 \(V=k^2\) 上由 \(s(a,b)=(-a,b)\) 作用。则
\[
k[V]^{C_2}=k[x^2,y],\qquad
\bigl(\operatorname{Sym}^2_k(V^*)\bigr)^{C_2}
=kx^2\oplus ky^2.
\]
\end{example}
\begin{proof}
由于 \(s^{-1}=s\)，对偶坐标满足 \(sx=-x,sy=y\)。将 \(f=\sum_{i,j}c_{ij}x^iy^j\) 代入 \(sf=f\)，逐个比较单项式系数，得到奇数 \(i\) 对应的 \(2c_{ij}=0\)，故这些系数全零。剩下的单项式都是 \((x^2)^r y^j\)，所以每个不变多项式都是 \(x^2,y\) 的多项式；反过来，\(x^2,y\) 都不变，它们的每个多项式也都不变。这给出了整个不变环，因为所分类的单项式覆盖任意次数。次数二时只剩 \(x^2,y^2\)。这个系数比较适用于有限域，无需用点集函数的相等代替多项式相等。
\end{proof}

\cref{b2:tab:08-group-actions}汇总本节出现的作用；同一张量幂上的两行分别来自 \(G\) 和 \(S_d\)。
\begin{table}[htbp]
\centering
\caption{表示及其张量构造上的群作用}\label{b2:tab:08-group-actions}
\begin{tabularx}{\linewidth}{@{}l l X@{}}
\toprule
空间 & 作用群 & 作用公式 \\
\midrule
\(V\) & \(G\) & \(v\mapsto \rho_V(g)v\) \\
\(V^*\) & \(G\) & \(\lambda\mapsto\lambda\circ\rho_V(g)^{-1}\) \\
\(V^{\otimes d}\) & \(G\) & \(\rho_V(g)^{\otimes d}(\bigotimes_i v_i)=\bigotimes_i\rho_V(g)v_i\) \\
\(V^{\otimes d}\) & \(S_d\) & \(P_\sigma(\bigotimes_i v_i)=\bigotimes_i v_{\sigma^{-1}(i)}\) \\
\(k[V]\) & \(G\) & \(g\lambda=\lambda\circ\rho_V(g)^{-1}\) 对 \(\lambda\in V^*\)，再按代数同态延拓 \\
\bottomrule
\end{tabularx}
\end{table}
\cref{b2:tab:08-group-actions}中的 \(G\) 在张量幂上同时变换各个因子，\(S_d\) 则只调换位置。下一命题先对整个 \(\operatorname{GL}_k(V)\) 证明两种作用相交换，因此也适用于通过 \(\rho_V\) 作用的任意群 \(G\)。

\subsection{张量表示与 Schur–Weyl 理论}
\label{b2:subsec:0804:03}

\begin{proposition}\label{b2:prop:commuting-tensor-actions}
设 \(k\) 为域，\(V\) 为 \(k\) 向量空间，\(d\ge1\)。在 \(V^{\otimes d}\) 上，\(\operatorname{GL}_k(V)\) 的对角作用 \(g\mapsto g^{\otimes d}\) 与 \(S_d\) 的位置置换作用 \(\sigma\mapsto P_\sigma\) 相交换：
\[
g^{\otimes d}P_\sigma=P_\sigma g^{\otimes d}.
\]
因此对称张量子空间和定义对称幂的关系子空间都在一般线性群作用下稳定。
\end{proposition}
\begin{proof}
对 \(v_1\otimes\cdots\otimes v_d\)，两个复合都得到
\[
gv_{\sigma^{-1}(1)}\otimes\cdots\otimes gv_{\sigma^{-1}(d)}.
\]
纯张量生成全空间，所以两算子相等。若 \(t\) 对所有置换不变，则 \(P_\sigma g^{\otimes d}t=g^{\otimes d}t\)；若一个关系为 \(t-P_\sigma t\)，其像为 \(g^{\otimes d}t-P_\sigma g^{\otimes d}t\)，仍是关系。于是两个稳定性结论分别成立。
\end{proof}

这已经说明两种对称性可以同时研究，但“作用交换”并不等于“每一种作用已经给出了另一种作用的全部对称算子”。为写清后一问题，设 \(k\) 为域，\(E\) 为 \(k\) 向量空间。若 \(\mathcal A\) 为 \(\operatorname{End}_k(E)\) 的子代数，记
\[
\operatorname{End}_{\mathcal A}(E)
=\bigl\{\,u\in\operatorname{End}_k(E)\mid ua=au\;\text{对每个}\;a\in\mathcal A\,\bigr\},
\]
称为该作用的\term[zhongxinhua-zi]{中心化子}[centralizer]。它是子代数，因为恒等、和、积都继续与 \(\mathcal A\) 交换。

\ProseParagraphBreak
也可以从一个算子族 \(\mathcal S\) 出发，令 \(\mathcal A\) 为它生成的含单位子代数。算子 \(u\) 若与每个 \(s\in\mathcal S\) 交换，就与这些算子的任意有限乘积及线性组合交换，因而与整个 \(\mathcal A\) 交换；反向由 \(\mathcal S\subseteq\mathcal A\) 立即成立。因此求中心化子时只须检验生成算子。对于张量幂上的两种作用，所要问的是：位置置换生成的代数是否已经包含所有与一般线性群作用交换的算子，反方向又是否成立？

\begin{claim}[Schur–Weyl 对偶]\label{b2:claim:schur-weyl-preview}
设 \(V\) 为有限维复向量空间，\(d\ge1\)。在 \(E=V^{\otimes d}\) 上，令 \(\mathcal A\) 为位置置换算子 \(P_\sigma\) 的复线性生成空间，\(\mathcal B\) 为算子 \(g^{\otimes d}\)、\(g\in\operatorname{GL}(V)\) 的复线性生成空间。此时
\[
\mathcal A=\operatorname{End}_{\mathcal B}(E),\qquad
\mathcal B=\operatorname{End}_{\mathcal A}(E).
\]
\end{claim}

两种线性生成空间确为含单位子代数，因为各自生成算子的乘积仍属于同类算子。\cref{b2:prop:commuting-tensor-actions}及算子线性组合的分配律只给出
\[
\mathcal A\subseteq\operatorname{End}_{\mathcal B}(E),\qquad
\mathcal B\subseteq\operatorname{End}_{\mathcal A}(E).
\]
两个反向包含才构成\term[Schur-Weyl-duiou]{Schur–Weyl 对偶}[Schur–Weyl duality]的实质：它们断言中心化子中没有其他算子。若 \(V=0\)，两等式的两侧均为零；非零情形的证明见第五卷第7.3节，也可参阅 Etingof 等人的《Introduction to Representation Theory》\cite[Theorem 5.18.4, p. 121; Proposition 5.19.1, p. 122]{EtingofEtAl2011RepresentationTheory}。

\ProseParagraphBreak
交换性本身不能推出这些等号。例如在 \(V=k^2\)、\(E=V^{\otimes2}\) 上，把一般线性群的作用换成平凡群的作用，其生成代数便只有 \(\mathcal B=k\operatorname{id}_E\)。位置置换代数为 \(\mathcal A=k\operatorname{id}_E+kP\)，其中 \(P\) 交换两个因子。两种作用当然交换，但 \(\operatorname{End}_{\mathcal B}(E)=\operatorname{End}_k(E)\) 的维数为十六，而 \(\mathcal A\) 的维数至多为二，所以第一个包含严格。

\ProseParagraphBreak
对域 \(k\) 上的向量空间 \(V\)，若 \(2\) 在 \(k\) 中可逆，令 \(P\) 为 \(V^{\otimes2}\) 上交换两个因子的算子，则两个算子
\[
E_+=\frac{\operatorname{id}+P}{2},\qquad
E_-=\frac{\operatorname{id}-P}{2}
\]
满足 \(E_++E_-=\operatorname{id}\)、\(E_\pm^2=E_\pm\)、\(E_+E_-=E_-E_+=0\)，因为 \(P^2=\operatorname{id}\)。这里 \(E_+\) 正是\cref{b2:prop:symmetrization-characteristic}在 \(S_2\) 上的平均算子，\(E_-\) 则是它的互补投影。因此
\[
V^{\otimes2}=\operatorname{im}E_+\oplus\operatorname{im}E_-,
\]
两项分别是 \(P\) 特征值为一的对称张量子空间和特征值为负一的反对称张量子空间；由\cref{b2:prop:commuting-tensor-actions}，它们在一般线性群作用下都稳定。下一章将通过外代数解释第二部分，并专门处理特征二时这两种条件不再分离的情形。


\begin{chaptersummary}
多重张量积由多线性映射的泛性质确定。结合和位置置换都由纯张量上的明确公式给出，生成性保证各种相容等式。张量代数把所有次数作直和并连接张量串，因而适合向任意结合代数指定线性生成元的像；只有底模具有一组自由基时，才得到以所有有限字为基的自由结合代数。

\ProseParagraphBreak
对称代数由交换关系取商，有限自由模的基使它成为多项式代数。增广理想模去其平方恢复一次生成模；这两种代数的构造也与扩张标量相容，无需平坦假设。对称代数的次数分量是对称商，张量幂中的对称张量则是不变子模；阶乘可逆时，平均算子给出两者的自然同构。即使在特征零的整数环上，不可逆的阶乘也会妨碍这个同构；正特征的轨道和还可能在商中消失，取幂也只满足相应的标量幂次规律。

\ProseParagraphBreak
群表示与张量运算相容，对偶作用中的逆元使求值配对保持不变，并保证复合次序正确。有限维条件把 \(V^*\otimes W\) 识别为 \(\operatorname{Hom}(V,W)\)，其不变量正是等变映射。多项式代数 \(k[V]=\operatorname{Sym}_k(V^*)\) 使用同一种对偶作用；其各次不动空间连同乘法组成不变环，有限域上仍须保留形式多项式与点集函数的区别。位置置换与一般线性群的作用交换；Schur–Weyl 对偶在有限维复空间上进一步给出双中心化子等式，其证明归于第五卷第7.3节。下一章改用交替关系，研究外幂与行列式。
\end{chaptersummary}

\BookExercises

\begin{exercise}\label{b2:ex:08-trilinear-coordinates}
设 \(M_1=R^2,M_2=R^3,M_3=R^2\)，\(N\) 为 \(R\) 模。给定十二个元素 \(a_{ijk}\in N\)。证明存在唯一三线性映射，使标准基向量组 \((e_i,f_j,e_k)\) 的像为 \(a_{ijk}\)，写出任意三元组的像。若 \(R=k\) 为域且 \(\dim_kN=s<\infty\)，求全部这些三线性映射组成的空间的维数。
\end{exercise}
\begin{solution}
多线性强迫
\[
b\left(\sum_i x_i e_i,\sum_j y_jf_j,\sum_k z_ke_k\right)
=\sum_{i,j,k}x_iy_jz_k a_{ijk}.
\]
各坐标唯一，因此公式良定义；对每一个变量分别核对加法与标量乘法，得到三线性，所以存在且唯一。由\cref{b2:prop:multiple-tensor-basis}，定义域张量积的基有十二个元素，其线性映射到 \(N\) 可以在这十二个元素上任意取值，故同构于 \(N^{12}\)，维数为 \(12s\)。
\end{solution}

\begin{exercise}\label{b2:ex:08-permutation-convention}
设 \(V\) 有三个不同基向量 \(e_1,e_2,e_3\)，取 \(\sigma=(12)\)、\(\pi=(23)\)。分别计算 \(P_\sigma P_\pi(e_1\otimes e_2\otimes e_3)\) 与 \(P_\pi P_\sigma(e_1\otimes e_2\otimes e_3)\)。若将定义改成按 \(\sigma(1),\ldots,\sigma(d)\) 排列，复合次序将如何变化？
\end{exercise}
\begin{solution}
先 \(P_\pi\) 后 \(P_\sigma\) 得到 \(e_3\otimes e_1\otimes e_2\)；先 \(P_\sigma\) 后 \(P_\pi\) 得到 \(e_2\otimes e_3\otimes e_1\)。两者不同，因为它们是不同的纯基张量。若改用 \(Q_\sigma(\bigotimes_jv_j)=\bigotimes_jv_{\sigma(j)}\)，则 \(Q_\sigma=P_{\sigma^{-1}}\)，故
\[
Q_\sigma Q_\pi=P_{\sigma^{-1}\pi^{-1}}
=P_{(\pi\sigma)^{-1}}=Q_{\pi\sigma}.
\]
于是得到反序的复合法则。置换作用的公式必须与采用的左、右作用约定一起说明。
\end{solution}

\begin{exercise}\label{b2:ex:08-multilinear-dual}
设 \(M_1,\ldots,M_d\) 为有限自由 \(R\) 模。证明
\[
M_1^*\otimes_R\cdots\otimes_RM_d^*
\longrightarrow\operatorname{Hom}_R(M_1\otimes_R\cdots\otimes_RM_d,R)
\]
由 \((\lambda_1\otimes\cdots\otimes\lambda_d)(m_1\otimes\cdots\otimes m_d)=\prod_j\lambda_j(m_j)\) 定义，是自然同构，其中 \(M_j^*=\operatorname{Hom}_R(M_j,R)\)。
\end{exercise}
\begin{solution}
先固定各 \(\lambda_j\)，右端对 \(m_j\) 多线性，故定义张量积上的泛函；所得泛函又对各 \(\lambda_j\) 多线性，所以诱导所写线性映射。对每个 \(M_j\) 选有限基及坐标对偶基。由\cref{b2:prop:multiple-tensor-basis}，纯基张量构成定义域基，而纯对偶基张量的像恰是目标中对应纯基张量的坐标泛函，故映射为同构。给定各 \(f_j:N_j\rightarrow M_j\)，两侧的预复合在纯张量上都取值为 \(\prod_j\lambda_j\bigl(f_j(n_j)\bigr)\)，证明自然性。
\end{solution}

\begin{exercise}\label{b2:ex:08-three-tensor-not-pure}
在 \((k^2)^{\otimes3}\) 中证明 \(t=e_1\otimes e_1\otimes e_1+e_2\otimes e_2\otimes e_2\) 不是纯张量。由此说明任意张量虽能写成纯张量的有限和，也不能总压缩成一项。
\end{exercise}
\begin{solution}
若 \(t=a\otimes b\otimes c\)，写 \(a=a_1e_1+a_2e_2\)，类似写 \(b,c\)。纯基张量的坐标唯一，给出 \(a_1b_1c_1=1\)、\(a_2b_2c_2=1\)，故六个系数都非零。但 \(e_1\otimes e_1\otimes e_2\) 的系数应为零，即 \(a_1b_1c_2=0\)，与域中无零因子矛盾。因此 \(t\) 不是纯张量，而题中已经给出两项纯张量之和。
\end{solution}

\begin{exercise}\label{b2:ex:08-tensor-torsion-module}
对整数 \(n>1\)，证明 \(T_{\mathbb Z}(\mathbb Z/n\mathbb Z)\cong\mathbb Z[X]/(nX)\) 作为分次代数成立，生成元 \(\bar1\) 对应 \(X\) 的像。指出每个次数的分量，以及它为何不是一个自由 \(\mathbb Z\) 模。
\end{exercise}
\begin{solution}
由\cref{b2:prop:tensor-quotient}，每个正次张量幂均为 \(\mathbb Z/n\mathbb Z\)，由 \(\bar1^{\otimes d}\) 生成；零次为 \(\mathbb Z\)。乘法把这两个纯幂连接为更高次纯幂。另一方面，理想 \((nX)\) 恰由常数项为零、所有正次系数均被 \(n\) 整除的多项式组成。因此商多项式的零次系数在 \(\mathbb Z\)，各正次系数在 \(\mathbb Z/n\mathbb Z\)，乘法正好相同。逐次数对应即给出分次同构。一次分量中的非零元素 \(\bar1\) 被 \(n\) 零化，故整个底模有挠；非零自由 \(\mathbb Z\) 模无挠，所以不是自由模。
\end{solution}

\begin{exercise}\label{b2:ex:08-augmentation}
设 \(R\) 为交换含幺环，\(M\) 为幺性 \(R\) 模，\(I\) 为 \(T_R(M)\) 的增广理想。对每个 \(r\ge1\)，求 \(I^r\) 的各次分量，证明 \(I^r/I^{r+1}\cong M^{\otimes r}\)，并证明这些同构使由乘法诱导的映射
\[
(I^r/I^{r+1})\otimes_R(I^s/I^{s+1})\longrightarrow I^{r+s}/I^{r+s+1}
\]
对应于连接张量串的同构，其中 \(r,s\ge1\)。再取 \(M=kx\)、\(R=k\) 为域，证明由 \(x\mapsto x+x^2\) 定义的增广代数自同态在 \(I/I^2\) 上诱导恒等，却不是代数自同构。
\end{exercise}
\begin{solution}
每个 \(I\) 因子只有正次分量，\(r\) 个这样的因子之积没有低于 \(r\) 次的分量，故 \(I^r\subseteq\bigoplus_{d\ge r}M^{\otimes d}\)。反过来，当 \(d\ge r\) 时，每个 \(d\) 重纯张量可以分成 \(r\) 个非空张量串的乘积，例如把前 \(r-1\) 个一次因子分别取出，其余合为最后一串。这些因子均在 \(I\) 中，而纯张量生成整个次数分量，故
\[
I^r=\bigoplus_{d\ge r}M^{\otimes d}.
\]
取次数 \(r\) 的分量，便给出 \(I^r/I^{r+1}\cong M^{\otimes r}\)。如果改变代表元，乘积之差落在 \(I^{r+s+1}\)，所以题述乘法在商上良定义。它在纯张量上连接两个串，由\cref{b2:thm:tensor-associativity-unit}得到所述同构。整个论证只用了纯张量的生成性，不要求 \(M\) 自由。

当 \(M=kx\) 时，\(T_k(M)=k[x]\)。赋值 \(x\mapsto x+x^2\) 给出保增广的自同态，且 \(x+x^2\equiv x\pmod{I^2}\)，所以一次商上的映射为恒等。但若 \(x\) 在此自同态的像中，就有多项式 \(f\) 满足 \(f(x+x^2)=x\)。\(f\) 不可能是常数；若其次数为 \(d\ge1\)，左侧的次数为 \(2d\)，不可能等于一。因此自同态不满射。一次商记录线性部分，并不足以判定整个代数映射是否可逆。
\end{solution}

\begin{exercise}\label{b2:ex:08-tensor-algebra-injection-failure}
令 \(R=k[t]\)，\(M=(t)\subset R\)，其中 \(k\) 为域。写出包含 \(i:M\hookrightarrow R\) 诱导的张量代数同态的生成元公式，证明它单射但不满射，并求其各次分量上的余核。再把系数环扩张到 \(k(t)\)，证明所得代数同态为同构。
\end{exercise}
\begin{solution}
\(M\) 是以 \(t\) 为基的秩一自由 \(R\) 模，而目标 \(R\) 以 \(1\) 为基。把它们在各自张量代数中的一次像分别记为 \(X,Y\)，由\cref{b2:prop:free-associative-algebra}，所求同态写成
\[
R[X]\longrightarrow R[Y],\qquad X\longmapsto tY.
\]
一次生成元 \(X\) 代表模中的基元素 \(t\)，不能把它与零次系数 \(t\) 混同。次数 \(d\) 的映射是 \(RX^d\to RY^d\)、\(aX^d\mapsto at^dY^d\)。\(R\) 为整环，故每个分量上的映射均单射，整个直和上的映射也单射。次数零的余核为零；\(d\ge1\) 时余核为 \(R/(t^d)\)。像是 \(R[tY]\)，而 \(Y\) 不在像中，因为它的一次系数 \(1\) 不属于 \(tR\)，所以映射不满射。

由\cref{b2:prop:tensor-symmetric-base-change}，扩张标量后得到 \(k(t)[X]\to k(t)[Y]\)、\(X\mapsto tY\)。此时 \(t\) 可逆，反向赋值 \(Y\mapsto t^{-1}X\) 给出逆同态。\secref{b2:sec:0802}的双数环反例说明张量代数不普遍保持单射；本题则说明必须检查具体底环、底模及各次张量映射，不能把该反例误读成每个非平凡包含都失去单射性。
\end{solution}

\begin{exercise}\label{b2:ex:08-inhomogeneous-relation}
设 \(k\) 为域，\(A=k[X]/(X^2-X)\)。证明 \(A\) 中 \(1,x\) 为一组基，其中 \(x\) 是 \(X\) 的像；证明这个商不能继承把 \(x\) 放在次数一的多项式分次。
\end{exercise}
\begin{solution}
由首一多项式的带余除法，每个商元素唯一表示为 \(a+bx\)，所以 \(1,x\) 是基，尤其 \(x\ne0\)。若继承原分次，则 \(x\) 在次数一、\(x^2\) 在次数二。关系 \(x^2=x\ne0\) 使同一个非零元素同时落入两个不同次数分量，与直和分解矛盾。商仍然是代数，失去的只是所指定次数的分次。
\end{solution}

\begin{exercise}\label{b2:ex:08-square-zero-algebra}
证明 \(k\langle x,y\rangle/(x^2,xy,yx,y^2)\) 有基 \(1,x,y\)。除了用关系约化说明生成性，还须证明线性无关性。
\end{exercise}
\begin{solution}
任意长度至少二的字都含一个所列二字关系，故在商中为零，得到生成性。要检测这三个元素是否独立，可以构造一个仍保留单位元及两个独立一次元素、而把所有二次乘积都置零的代数。在向量空间 \(B=k\oplus k^2\) 上定义
\[
(a,u)(b,v)=(ab,av+bu).
\]
单位为 \((1,0)\)。两种三项乘积的第二坐标都等于 \(ab w+ac v+bc u\)，第一坐标都为 \(abc\)，故结合；它是 \(k\) 代数。令 \(x\mapsto(0,e_1)\)、\(y\mapsto(0,e_2)\)，所有二字关系的像为零，故\cref{b2:prop:algebra-relations-universal}给出商到 \(B\) 的同态。\(1,x,y\) 的像是 \(B\) 的一组基，所以它们在原商中线性无关。
\end{solution}

\begin{exercise}\label{b2:ex:08-symmetric-quotient}
设 \(N\subseteq M\) 为子模。令 \((N)\) 表示 \(\operatorname{Sym}_R(M)\) 中由 \(N\) 的一次像生成的理想。证明
\[
\operatorname{Sym}_R(M/N)\cong\operatorname{Sym}_R(M)/(N)
\]
为自然分次代数同构。
\end{exercise}
\begin{solution}
商映射 \(M\rightarrow M/N\) 诱导对称代数同态，并将 \(N\) 送到零，故得到右端商代数到左端的同态。反过来，映射 \(M\rightarrow\operatorname{Sym}_R(M)/(N)\) 在 \(N\) 上为零，故公式 \(m+N\mapsto\bar m\) 良定义且线性。由\cref{b2:thm:symmetric-algebra-universal}，它延拓为左端到右端的代数同态。两复合在一次生成元和结构标量上均为恒等，所以互逆。构造保持一次分量和次数，也与保持子模的线性映射相容，故为自然分次同构。
\end{solution}

\begin{exercise}\label{b2:ex:08-symmetric-direct-sum}
对交换 \(R\) 代数 \(A,B\)，先核对 \(A\otimes_RB\) 由 \((a\otimes b)(a'\otimes b')=aa'\otimes bb'\) 成为交换 \(R\) 代数。再证明
\[
\operatorname{Sym}_R(M\oplus N)\cong
\operatorname{Sym}_R(M)\otimes_R\operatorname{Sym}_R(N),
\]
使 \((m,0)\) 对应 \(m\otimes1\)、\((0,n)\) 对应 \(1\otimes n\)。
\end{exercise}
\begin{solution}
四变量公式 \(aa'\otimes bb'\) 对每个变量线性，且把两对因子中的 \(R\) 平衡关系分别送成相等元素，因此诱导张量积上的双线性乘法。在纯张量上核对结合、交换及单位 \(1\otimes1\)，再由生成性推广，便得交换代数。

题述一次映射由 \(M\oplus N\) 的直和性质定义，继而由对称代数泛性质延拓为正向同态。两次包含 \(M,N\rightarrow M\oplus N\) 分别诱导同态 \(\alpha,\beta\) 到左端；公式 \(a\otimes b\mapsto\alpha(a)\beta(b)\) 平衡，且目标交换，故诱导反向代数同态。两复合在 \(M,N\) 的一次像上为恒等；右端由 \(a\otimes1\)、\(1\otimes b\) 生成，而这些又由相应一次像及标量生成，因此两个复合在全代数为恒等。
\end{solution}

\begin{exercise}\label{b2:ex:08-symmetric-dimension}
若 \(\dim_kV=3\)，求 \(V^{\otimes4}\)、\(\operatorname{Sym}_k^4(V)\) 和商映射 \(q_4\) 的核的维数。说明这些数是否依赖域的特征，并给出核中的一个非零元素。
\end{exercise}
\begin{solution}
张量幂维数为 \(3^4=81\)，对称幂维数为 \(\binom64=15\)。\(q_4\) 满射，故核维数为 \(81-15=66\)。前两项的基分别由有序下标与全次数四的单项式给出，都不涉及除法，因此任何特征下相同。取两个不同基向量 \(e_1,e_2\)，则
\[
e_1\otimes e_2\otimes e_1\otimes e_1
-e_2\otimes e_1\otimes e_1\otimes e_1
\]
是两个不同纯基张量之差，非零且在对称商中为零，特征二时也仍非零。
\end{solution}

\begin{exercise}\label{b2:ex:08-symmetrization-orbits}
设 \(k\) 为域，\(V=ke\oplus kf\)。求 \((V^{\otimes4})^{S_4}\) 的一组基，并写出商映射在这组基上相对于 \(e^4,e^3f,e^2f^2,ef^3,f^4\) 的矩阵。分别计算 \(k=\mathbb Q\)、特征二与特征三时的核和余核，并在 \(\mathbb Q\) 上写出 \(e^2f^2\) 对应的平均张量。若把底环改为 \(\mathbb Z\)，求同一限制商映射的核和余核。
\end{exercise}
\begin{solution}
十六个纯基张量按 \(f\) 出现的次数 \(j=0,1,2,3,4\) 分成五个轨道。记 \(u_j\) 为含 \(j\) 个 \(f\) 的全部不同纯基张量之和。不变性要求同一轨道的系数相同，各轨道的支集互不相交，故 \(u_0,\ldots,u_4\) 是不变部分的一组基。第 \(j\) 个轨道有 \(\binom4j\) 项，每项的商像均为 \(e^{4-j}f^j\)，所以矩阵为
\[
\operatorname{diag}(1,4,6,4,1).
\]
在 \(\mathbb Q\) 上这五个数可逆，核与余核均为零。\(e^2f^2\) 对应 \(u_2/6\)：在二十四项平均中，每个不同排列重复 \(2!2!=4\) 次，所以其系数为 \(4/24=1/6\)。

特征二时矩阵为 \(\operatorname{diag}(1,0,0,0,1)\)，核以 \(u_1,u_2,u_3\) 为基，余核以 \(e^3f,e^2f^2,ef^3\) 的类为基，均为三维。特征三时矩阵为 \(\operatorname{diag}(1,1,0,1,1)\)，核为 \(ku_2\)，余核由 \(e^2f^2\) 的类生成，均为一维。这与\cref{b2:exm:cubic-symmetrization-orbits}中三次、特征二时的同构不同：不能只由阶乘不可逆判定某个具体次数的映射。

在整数环上，同一轨道系数的论证仍成立，矩阵的五个对角元都非零，故核为零；逐坐标取商得到余核
\[
\mathbb Z/4\mathbb Z\oplus\mathbb Z/6\mathbb Z\oplus\mathbb Z/4\mathbb Z.
\]
\end{solution}

\begin{exercise}\label{b2:ex:08-invariants-over-integers}
令 \(M=\mathbb Z^r\)、\(r\ge2\)，标准基为 \(e_1,\ldots,e_r\)。求 \((M^{\otimes2})^{S_2}\) 的一组基，并写出限制商映射到 \(\operatorname{Sym}_{\mathbb Z}^2(M)\) 的矩阵。求其核和余核，再求把此矩阵模二后所得线性映射的核与余核的维数。
\end{exercise}
\begin{solution}
交换因子固定每个 \(e_i\otimes e_i\)，并交换 \(e_i\otimes e_j\) 与 \(e_j\otimes e_i\)。不变性要求每对交换项的系数相同，所以不变量以
\[
e_i\otimes e_i\quad(1\le i\le r),\qquad
e_i\otimes e_j+e_j\otimes e_i\quad(1\le i<j\le r)
\]
为基。目标依次以 \(e_i^2\) 和 \(e_ie_j\) 为基时，矩阵为
\[
\operatorname{diag}(\underbrace{1,\ldots,1}_{r\text{ 项}},
\underbrace{2,\ldots,2}_{\binom r2\text{ 项}}).
\]
所有对角元非零，限制映射单射，余核为 \((\mathbb Z/2\mathbb Z)^{\binom r2}\)，由不同 \(e_ie_j\) 的类给出独立的二阶元素。模二后，对角项 \(1\) 保留、\(2\) 消失，故核与余核的维数均为 \(\binom r2\)。章首的秩二计算是这一计算的最小情形；增加基向量后，每一对不同下标都贡献一个不能在整数环上平均的混合项。
\end{solution}

\begin{exercise}\label{b2:ex:08-frobenius-semilinear}
在 \(k=\mathbb F_2[a]/(a^2+a+1)\) 上，令 \(V=ke\)，并定义 \(\Phi_2:V\rightarrow\operatorname{Sym}_k^2(V)\)、\(v\mapsto v^2\)。求 Frobenius 同态 \(F_k:c\mapsto c^2\) 在 \(a\) 上的值，证明 \(\Phi_2\) 可加且对 \(F_k\) 半线性，但不是 \(k\) 线性。若将底域换成 \(\mathbb F_2\)，同一公式有何变化？
\end{exercise}
\begin{solution}
多项式 \(a^2+a+1\) 在 \(\mathbb F_2\) 的两个元素 \(0\) 与 \(1\) 处取值均非零，所以商是四元域，且 \(F_k(a)=a^2=a+1\ne a\)。对称代数交换，特征二下交叉项消失，故 \(\Phi_2(u+v)=\Phi_2(u)+\Phi_2(v)\)；对任意 \(c\in k\) 又有 \(\Phi_2(cv)=c^2\Phi_2(v)=F_k(c)\Phi_2(v)\)。但
\[
(ae)^2=a^2e^2\ne ae^2,
\]
其中 \(e^2\ne0\) 由对称幂的单项式基保证，所以不是 \(k\) 线性。在 \(\mathbb F_2\) 上，每个标量满足 \(c^2=c\)，公式同时保持加法与标量乘法，因而成为线性映射。
\end{solution}

\begin{exercise}\label{b2:ex:08-polynomial-versus-function}
在 \(\mathbb F_3[X,Y]\) 中，证明齐次多项式 \(X^3Y-XY^3\) 非零，但在 \(\mathbb F_3^2\) 的每一点都取零。说明这对把 \(\operatorname{Sym}_{\mathbb F_3}^4\bigl((\mathbb F_3^2)^*\bigr)\) 当作函数空间有什么限制。
\end{exercise}
\begin{solution}
\(X^3Y\) 与 \(XY^3\) 是不同的次数四单项式，系数分别为一和负一，故形式多项式非零。对 \(\mathbb F_3\) 中的每个元素都有 \(a^3=a\)，所以在 \((a,b)\) 处取值为 \(ab-ab=0\)。选取对偶基后，对称幂与这些形式齐次多项式同构，所给非零元素却成为零函数。因此取值映射不单射，不能用函数相等代替对称代数中的元素相等。
\end{solution}

\begin{exercise}\label{b2:ex:08-traces}
设 \(A:V\rightarrow V\)、\(B:W\rightarrow W\) 为有限维空间的线性算子。证明
\(\operatorname{tr}(A\otimes B)=\operatorname{tr}(A)\operatorname{tr}(B)\)。
若 \(2\ne0\) 于 \(k\)，进一步证明
\[
\operatorname{tr}\bigl(\operatorname{Sym}^2(A)\bigr)
=\frac{(\operatorname{tr}A)^2+\operatorname{tr}(A^2)}2.
\]
\end{exercise}
\begin{solution}
在张量基 \(e_i\otimes f_j\) 上，\(A\otimes B\) 的对应对角元为 \(a_{ii}b_{jj}\)，求和得第一式。令 \(P\) 为交换两因子的算子。\(A\otimes A\) 与 \(P\) 交换，因而保持 \(P\) 的两个特征子空间。由\cref{b2:prop:symmetrization-characteristic}，投影 \(E_+=(1+P)/2\) 的像自然识别为对称商。在对称子空间上，\(E_+(A\otimes A)\) 等于 \(\operatorname{Sym}^2(A)\)；在反对称子空间上它为零。两部分的直和就是整个张量空间，所以整个空间上的迹等于对称部分上的迹，即
\[
\operatorname{tr}\bigl(\operatorname{Sym}^2(A)\bigr)
=\operatorname{tr}\bigl(E_+(A\otimes A)\bigr).
\]
\(P(A\otimes A)\) 在 \(e_i\otimes e_j\) 处的对角元为 \(a_{ji}a_{ij}\)，其总和为 \(\operatorname{tr}(A^2)\)。结合第一式并取平均，得到结论，无需假定 \(A\) 可对角化。
\end{solution}

\begin{exercise}\label{b2:ex:08-symmetric-unipotent-matrix}
设 \(A(e)=e\)、\(A(f)=e+f\)。写出 \(\operatorname{Sym}^2(A)\) 在 \(e^2,ef,f^2\) 上的矩阵，求特征二时与特征不为二时的固定子空间维数。
\end{exercise}
\begin{solution}
三列分别由
\[
A(e)^2=e^2,\quad A(e)A(f)=e^2+ef,\quad
A(f)^2=e^2+2ef+f^2
\]
给出，所以矩阵为 \(\begin{pmatrix}1&1&1\\0&1&2\\0&0&1\end{pmatrix}\)。固定向量的坐标 \((a,b,c)\) 满足 \(b+c=0\)、\(2c=0\)。特征不为二时 \(b=c=0\)，固定空间为 \(ke^2\)，维数一；特征二时只剩 \(b=c\)，固定空间由 \(e^2,ef+f^2\) 生成，维数二。对称幂的维数均为三，改变的是诱导算子的固定空间。
\end{solution}

\begin{exercise}\label{b2:ex:08-hom-invariant-jordan}
在 \(V=\mathbb Qe_1\oplus\mathbb Qe_2\) 上，让无穷循环群的生成元作用为
\(\begin{pmatrix}1&1\\0&1\end{pmatrix}\)。
求 \(\operatorname{End}_G(V)\)，并通过 \(\Theta\) 写出 \((V^*\otimes V)^G\) 的一组基。
\end{exercise}
\begin{solution}
设 \(C=\begin{pmatrix}a&b\\c&d\end{pmatrix}\)。与生成元矩阵相交换等价于 \(c=0,d=a\)，所以等变自同态恰为 \(aI+bE_{12}\)，构成二维空间。与生成元相交换也与其正负整数幂相交换，因此已核对全部群元素。由\cref{b2:prop:hom-as-tensor-representation}，\(\Theta\) 将不变张量对应到等变自同态，故所求不变张量的一组基是
\[
e_1^*\otimes e_1+e_2^*\otimes e_2,\qquad e_2^*\otimes e_1.
\]
它们分别对应 \(I\) 和 \(E_{12}\)。
\end{solution}

\begin{exercise}\label{b2:ex:08-commuting-not-centralizers}
设 \(V=\mathbb Q^2\)，\(E=V^{\otimes2}\)，\(P\) 为交换两因子的算子，\(\mathcal A=\operatorname{span}_{\mathbb Q}\{\operatorname{id}_E,P\}\)。利用对称与反对称子空间分解，求 \(\operatorname{End}_{\mathcal A}(E)\) 及其中心化子，证明后者恰为 \(\mathcal A\)。若令平凡群作用在 \(E\) 上，其生成代数为 \(\mathcal B=\mathbb Q\operatorname{id}_E\)，指出 \(\mathcal B\subseteq\operatorname{End}_{\mathcal A}(E)\) 是否为等号。
\end{exercise}
\begin{solution}
记 \(E_+=\ker(P-\operatorname{id})\)、\(E_-=\ker(P+\operatorname{id})\)。若 \(e_1,e_2\) 是 \(V\) 的标准基，则 \(E_+\) 有基
\[
e_1\otimes e_1,\quad e_1\otimes e_2+e_2\otimes e_1,\quad e_2\otimes e_2,
\]
而 \(E_-\) 以 \(e_1\otimes e_2-e_2\otimes e_1\) 为基。与 \(P\) 交换的算子恰为分别保持这两个特征子空间的算子，所以
\[
\mathcal C\coloneqq\operatorname{End}_{\mathcal A}(E)
\cong M_3(\mathbb Q)\oplus\mathbb Q,
\]
其维数为十。两块上的投影均在 \(\mathcal C\) 中；与 \(\mathcal C\) 交换的算子因而也须分块。其三维块须与全部三阶矩阵交换：与对角矩阵单位交换迫使该块为对角矩阵，与非对角矩阵单位交换又迫使三个对角元相等。故
\[
\operatorname{End}_{\mathcal C}(E)
=\{\,a\operatorname{id}_{E_+}\oplus b\operatorname{id}_{E_-}\mid a,b\in\mathbb Q\,\}
=\mathcal A,
\]
最后一个等号由 \(P\) 在两块上分别取 \(1,-1\) 得出。另一方面，\(\mathcal B\) 只有一维，严格小于十维的 \(\mathcal C\)。双中心化子的计算确实恢复了 \(\mathcal A\)，但与它交换的一个任意群作用，并不必生成整个 \(\mathcal C\)。
\end{solution}

\begin{exercise}\label{b2:ex:08-polynomial-weight-invariants}
令无穷循环群 \(G=\langle g\rangle\) 在 \(V=\mathbb Q^2\) 上由 \(g(a,b)=(2a,b/2)\) 作用。求对偶坐标 \(x,y\) 上的作用，证明 \(\mathbb Q[V]^G=\mathbb Q[xy]\)，并求每个次数 \(d\) 的齐次不变空间维数。
\end{exercise}
\begin{solution}
对偶作用使用 \(g^{-1}\)，故 \(gx=x/2,gy=2y\)。单项式 \(x^iy^j\) 的特征值是 \(2^{j-i}\)，在有理数域上恰当 \(i=j\) 时为一。逐项比较系数可知不变多项式正是有限和 \(\sum_i c_i(xy)^i\)。所以次数 \(d\) 为偶数时，不变空间由 \((xy)^{d/2}\) 生成，维数为一；次数为奇数时空间为零。这同时说明生成元 \(xy\) 在原分次下的次数为二。
\end{solution}

\begin{exercise}\label{b2:ex:08-algebra-base-change}
设 \(k\) 为域，\(R=k[\varepsilon]/(\varepsilon^2)\)、\(M=(\varepsilon)\subset R\)、\(S=R/(\varepsilon)=k\)。求 \(T_R(M)\) 与 \(\operatorname{Sym}_R(M)\) 的生成元关系表达，并分别计算它们扩张标量到 \(S\) 后的代数。证明 \(S\otimes_RM\) 非零，再计算包含 \(i:M\hookrightarrow R\) 在扩张标量后诱导的张量代数同态
\[
S\otimes_RT_R(M)\longrightarrow S\otimes_RT_R(R)
\]
的核及像。
\end{exercise}
\begin{solution}
模 \(M\) 由 \(\varepsilon\) 生成。写 \(a=a_0+a_1\varepsilon\in R\)，便有 \(a\varepsilon=a_0\varepsilon\)，故 \(\operatorname{Ann}_R(\varepsilon)=(\varepsilon)\)。因此它的全部线性关系由 \(\varepsilon\cdot\varepsilon=0\) 生成，\(M\cong R/(\varepsilon)\)。把这一模生成元记为代数中的 \(X\)，则
\[
T_R(M)\cong R[X]/(\varepsilon X).
\]
可用泛性质核对：到任意结合 \(R\) 代数 \(A\) 的线性赋值由一个满足 \(\varepsilon a=0\) 的元素 \(a\) 决定；一个生成元的自由结合代数就是 \(R[X]\)，该关系恰好允许赋值下降。这个代数已经交换，故 \(\operatorname{Sym}_R(M)\) 也为同一个商。

由\cref{b2:prop:tensor-quotient}，\(S\otimes_RM\cong M/\varepsilon M=M\)，作为 \(k\) 空间是一维，生成元 \(1\otimes\varepsilon\) 非零。\cref{b2:prop:tensor-symmetric-base-change}于是给出
\[
S\otimes_RT_R(M)\cong k[X],\qquad
S\otimes_R\operatorname{Sym}_R(M)\cong k[X].
\]
原来对一次生成元的 \(\varepsilon\) 关系，在新系数域中已自动成立；零次的 \(\varepsilon\) 也被商去。

另一方面，\(S\otimes_RT_R(R)\cong T_k(k)=k[Y]\)。扩张标量后的包含把 \(1\otimes\varepsilon\) 送到 \(S\otimes_RR\cong S\) 中的零，故诱导代数同态为 \(k[X]\to k[Y]\)、\(X\mapsto0\)。它把多项式送到常数项，核为 \((X)\)，像为常数子代数 \(k\)。扩张标量相容的同构在这里仍成立；失去单射性的是由 \(i\) 诱导的映射。
\end{solution}
